\section{Curriculum assurance: what is supported and what is not}
\label{curriculum-assurance}
\subsection{Conclusion and evidence standard}
The cyber tracks are \textbf{evidence-informed draft curricula with explicit
performance assessments}. They are not empirically proven courses, accredited
certifications, or a guarantee of career outcomes. Structural validation proves
that the software can parse and traverse a path; it does not prove that following
it produces intuition. Real evidence of an individual's learning is the ability
to explain a mechanism, reproduce a result, handle an unfamiliar variation, and
repeat the reasoning after a delay. A credible claim about the curriculum's
average effectiveness would require learner outcome data that do not yet exist.

This audit reviewed public first-party curricula, official exercise specifications,
research abstracts/publisher summaries, a teaching practice guide, and practitioner
accounts. It did not purchase courses, inspect private learner records, execute
every third-party lab, or run a controlled learning study. Source-specific access
and inference limits appear in the evidence register below. Published curricula
are useful alignment references, not external endorsements of Ascend.

\subsection{Research-to-design reasoning}
Roediger and Karpicke's 2006 experiments found an advantage for retrieval over
restudy on delayed tests of prose; Karpicke and Blunt's 2011 work also examined
comprehension/inference from science material. These are reasons to include
closed-notes reconstruction, not evidence that a memory quiz alone teaches
security operations. The Dunlosky review's publisher summary gives stronger
support to practice testing and distributed practice than passive rereading.
The IES guide recommends spacing, alternating worked examples with problem
solving, connecting concrete and abstract representations, and explanatory
questions. The resulting design combines explanation with actual lab behavior.

The practical sequence is: prerequisite demonstration; guided experiment;
prediction and observation; explanation of mechanism; changed-case assessment;
positive/negative control; later retrieval and correction. References and hints
are legitimate during learning. A learner who needed a complete walkthrough
records that help, practices, then attempts a fresh small variation before claiming
independent transfer. There is no rule that struggling without guidance is better
than learning from a worked example.

\subsection{Lessons from practitioners and their limits}
Julia Evans describes identifying the exact question behind vague confusion and
researching it until her model improves. LiveOverflow describes programming,
CTF/wargame attempts, and studying unsolved challenge write-ups as parts of his
own development. The useful adaptation is a cycle of attempt, feedback and a new
attempt, not his page's unsupported universal 10,000-hour slogan.

Vendor-hosted interviews with Kamil Vavra and Johnny Villarreal give concrete
accounts of lab practice informing workplace understanding, while also describing
help from peers and difficulty with unfamiliar mechanisms. These are selected,
self-reported successes subject to survivorship and marketing bias. They do not
justify guarantees, all-night work, compulsive platform use, or the claim that one
person's ordering suits every beginner. James Kettle's published research is a
technical example of testing a mechanism across implementations, not a curriculum
outcome study. Ascend uses these accounts as context, below experimental learning
research in the evidence hierarchy.

\subsection{Scope and prerequisites of each track}
\begin{longtable}{p{.22\textwidth}p{.32\textwidth}p{.38\textwidth}}
\textbf{Track} & \textbf{Entry demonstration} & \textbf{Outcome and limits}\\\hline
Networking & Local service, address/port, client request, capture location. & Packet/route reasoning and bounded troubleshooting. Does not cover all enterprise routing, IPv6 engineering, or wireless exploitation.\\
Web application security & Capture/replay a training request and explain input, response and identity. & Selected web vulnerability mechanisms, authorization testing, and repair. Not the entire OWASP taxonomy or a replacement for the full Academy.\\
Network penetration testing & Owned/training host, shell identity, network baseline and reset. A separate Windows baseline precedes Windows escalation. & Scoped enumeration, service hypotheses, lab escalation and reporting. Not the full enterprise/job-role syllabus, stealth training, or authority to test third parties.\\
Reverse engineering & Compile simple C; explain arrays, pointers, lifetime and ABI. & Small-binary reasoning and introductory memory-corruption/ROP mechanisms. Not a guarantee of kernel/browser exploit expertise. Architecture and lab image must match.\\
Defensive security & Available telemetry, user permissions, time handling and a harmless recorded event. & Log investigation, bounded detection evaluation, incident evidence and controls. Enterprise Windows DFIR and other SOC specialisms need additional training.\\
\end{longtable}
Foundation failures are useful diagnostic results. Use the linked foundation
material and repeat a small demonstration; do not treat a tool installation as
proof of the underlying prerequisite. Main milestones depend on the level's
preparatory tasks, including previously omitted patching and transfer concepts.

\subsection{Technical corrections made by the audit}
Phoenix Stack Zero overwrites adjacent data; later exercises introduce indirect
calls or saved-return-address control. ROP Emporium's multi-call assessment now
names \emph{callme}, rather than implying that every split/ret2win solution shows
multi-call chaining. Examples record architecture, ABI, allocator and mitigations.
Undefined behavior need not crash, and a decompiler need not make an error for an
exercise to be useful.

Security Onion only supports x86-64. A Pi may supply telemetry, but cannot host
that distribution. Import/Eval/Standalone are different modes. The current
hardware table says 50 GB for Import while the first-time guide says 200 GB;
plan conservatively for 200 GB and check the chosen release. A supported hosted
SIEM lab is an alternative. No subscription or infrastructure cost is silently
assumed to be zero. Import exercises teach ingestion/search, not continuous
production monitoring.

A default Suricata ruleset does not guarantee an alert for an arbitrary nmap scan.
A DHCP renewal need not contain all four DORA messages. An address lookup does not
by itself show NAT rewriting; a public address comparison is not inside/outside
packet evidence. GraphQL introspection alone is not an authorization vulnerability.
Juice Shop is not assumed to provide a standard access log without configuration.
Sigma event selection is distinguished from backend-dependent correlation.
No exercise requires a learner to invent an error, vulnerability, real outage,
or unobservable fact in order to complete it.

The final version-3 resource audit checked 117 unique URLs: 111 responded
successfully and six returned access restrictions (TryHackMe 429 or Cloudflare
403). Unreachable CTF101 references and failing Exploit-DB references were replaced
with relevant official Phoenix, pwn.college, ROP Emporium and Rapid7 resources.
The replacements were opened and inspected; a reachable index still requires
selecting the relevant exercise. The machine-readable audit in
\path{docs/testing/cyber-track-links-v3.json} preserves original failures and
replacements. HTTP success is neither a full content audit nor a promise that a
paid or account-gated lab is freely available. Restricted responses are not proof
that the resource is dead.

\subsection{Assessment rubric and evidence packet}
For each milestone, retain the lab/dataset version and scope, baseline prediction,
observed evidence, reproduction steps, a changed case, a control, the explanation,
limitations and help used. Keep third-party solutions/flags private where required
and exclude credentials. The following are separate decisions, not a score whose
high average can hide a missing skill:
\begin{enumerate}
\item Can the learner explain the causal mechanism rather than only recite a command?
\item Can the result be reproduced from a known state using the saved evidence?
\item Can the learner independently solve or correctly predict a changed case?
\item Does a contrasting control distinguish the claimed mechanism from noise,
configuration failure, or a broken legitimate workflow?
\end{enumerate}
Claim the assessment only when all four are supported. If help reveals the answer,
record practice and attempt a fresh variation. A peer or instructor should select
or inspect the variant when possible; self-selected easy variants are weaker
evidence. The existing file verifier only checks artifact presence and optional
size/pattern. It is not an automated grader of these four decisions.

Each level also has an optional return-later node. Attempt it on another day,
usually 2--7 days later, before opening the original notes; then compare, correct,
and log what remains uncertain. This interval is an authoring choice, not a
validated optimal schedule. Ascend currently checks only the prerequisite,
not elapsed time. Reopening a practice task is possible; there is no automated
spaced-repetition scheduler or retention certificate.

\subsection{Versioning and existing learners}
Version 3 preserves all version-2 task IDs and prior progress. It adds five entry
checks, one Windows baseline and forty return-later tasks. A previously completed
milestone remains completed for navigation, but its old completion does not
retroactively meet the new assessment criteria. Existing learners can use the
new recall task or explicitly reopen a milestone to collect evidence under this
rubric. Review/trust resets on changed pack content; neither this audit nor a
successful file check grants human review status. Every pack remains an AI draft.

\subsection{How to evaluate learning in actual use}
Before a level, attempt a short baseline task without its solution and record
explanation, result, confidence, time and help. After its guided work, attempt a
fresh comparable task using the four assessment decisions above. Repeat a small
changed-case test after about a week. Track which misconceptions disappear, which
skills transfer, assistance required, and whether gains survive the delay.
Completion count, confidence, streak length and minutes alone are not learning
outcome measures. Failure should lead to a smaller experiment or prerequisite
repair, not invented success evidence.

For a pilot with several learners, collect starting experience, hardware/tool
constraints, baseline and delayed task results, attrition and time costs. Use a
predefined rubric and, if feasible, an assessor unaware of whether an answer was
pre- or post-training. Report missing follow-ups and failed tasks. A before/after
pilot can support feasibility and individual progress; causal comparisons with
another curriculum require a stronger study design. No such pilot has been run.

\subsection{Level-by-level assessment crosswalk}
Each entry corresponds to a named milestone in the original five or nineteen specialist tracks. Source IDs refer to the evidence register. The task text and prerequisites remain in the runtime TOML packs.
\subsubsection{networking-for-security: models}
\noindent\textbf{Milestone:} layers-milestone\par
\noindent\textbf{Skill:} Distinguish link, network, transport, and application evidence.\par
\noindent\textbf{Experiment:} Capture an owned HTTP exchange; identify addresses, ports, request and response by frame number. State what the capture does not reveal.\par
\noindent\textbf{Transfer:} Repeat with a different destination or port and predict which fields change.\par
\noindent\textbf{Control:} Compare a loopback capture with a LAN frame; do not invent Ethernet or physical-layer evidence on loopback.\par
\noindent\textbf{Later retrieval:} Sketch the layers from memory, then label a different frame.\par
\noindent\textbf{Alignment references:} htb-foundations, ies-practice\par
\subsubsection{networking-for-security: addressing}
\noindent\textbf{Milestone:} addressing-milestone\par
\noindent\textbf{Skill:} Predict subnet membership and next-hop selection.\par
\noindent\textbf{Experiment:} Design two small lab subnets and calculate network, host range, broadcast, and gateway; validate the arithmetic with Python ipaddress and inspect a matching routing table.\par
\noindent\textbf{Transfer:} Change the prefix or destination; predict whether delivery is local or via a gateway before checking.\par
\noindent\textbf{Control:} Include one off-subnet destination and one overlapping-subnet plan that must be rejected.\par
\noindent\textbf{Later retrieval:} Solve a fresh CIDR example and explain longest-prefix matching.\par
\noindent\textbf{Alignment references:} htb-foundations, ies-practice\par
\subsubsection{networking-for-security: transport}
\noindent\textbf{Milestone:} transport-milestone\par
\noindent\textbf{Skill:} Build a TCP exchange and reason about stream boundaries and connection failure.\par
\noindent\textbf{Experiment:} Build a local TCP client/server that exchanges a framed message. Demonstrate a successful exchange and a clean shutdown.\par
\noindent\textbf{Transfer:} Send the same message in several chunks and handle a response longer than one recv call.\par
\noindent\textbf{Control:} Stop the server and distinguish connection refusal from a server accepting but not replying; apply a timeout.\par
\noindent\textbf{Later retrieval:} Explain why TCP does not preserve application message boundaries, then repair one receive-loop error.\par
\noindent\textbf{Alignment references:} htb-foundations, ies-practice\par
\subsubsection{networking-for-security: application}
\noindent\textbf{Milestone:} application-milestone\par
\noindent\textbf{Skill:} Trace name resolution, connection establishment, and HTTP without assuming every page load is identical.\par
\noindent\textbf{Experiment:} Capture or correlate a controlled DNS lookup and an HTTP request to an owned lab service; identify each observed protocol.\par
\noindent\textbf{Transfer:} Compare a cached lookup or reused connection with the first request and predict which exchange may disappear.\par
\noindent\textbf{Control:} Compare an unresolvable name with a reachable host returning HTTP 404; separate DNS failure from application failure.\par
\noindent\textbf{Later retrieval:} Explain what HTTPS hides and what can remain visible; check a fresh request.\par
\noindent\textbf{Alignment references:} htb-foundations, ies-practice\par
\subsubsection{networking-for-security: capture}
\noindent\textbf{Milestone:} capture-milestone\par
\noindent\textbf{Skill:} Recover a conversation using packet filters and protocol documentation.\par
\noindent\textbf{Experiment:} Analyze a previously unused training capture or owned app exchange; identify a protocol and three fields with frame numbers. Standards and field references are allowed.\par
\noindent\textbf{Transfer:} Answer a new question from the same capture using a different filter or tshark extraction.\par
\noindent\textbf{Control:} Show a broad filter that includes unrelated traffic and explain how you excluded it.\par
\noindent\textbf{Later retrieval:} Write a filter from memory, then verify it on a different conversation.\par
\noindent\textbf{Alignment references:} htb-foundations, ies-practice\par
\subsubsection{networking-for-security: services}
\noindent\textbf{Milestone:} services-milestone\par
\noindent\textbf{Skill:} Enforce and verify a source-specific network policy without confusing it with an application outage.\par
\noindent\textbf{Experiment:} In a disposable lab with console access, allow one source to a test service and deny a second source. Save rules and both connection results.\par
\noindent\textbf{Transfer:} Move the test service to another port and predict the required rule change.\par
\noindent\textbf{Control:} Verify the service remains healthy from the allowed source and that restoring the original rules restores the baseline.\par
\noindent\textbf{Later retrieval:} Explain rule order and default policy, then predict the result of a small rule set.\par
\noindent\textbf{Alignment references:} htb-foundations, ies-practice\par
\subsubsection{networking-for-security: wireless}
\noindent\textbf{Milestone:} wireless-milestone\par
\noindent\textbf{Skill:} Assess Wi-Fi authentication and guest isolation using observed configuration and connectivity.\par
\noindent\textbf{Experiment:} Audit your owned access point or a provided lab configuration: authentication mode, WPS, guest isolation, and administrative exposure. Test guest-to-LAN access where supported.\par
\noindent\textbf{Transfer:} Compare two configurations and predict which permits guest access to a chosen LAN service.\par
\noindent\textbf{Control:} Check that an isolated guest can still reach an allowed destination; a broken guest connection does not prove isolation.\par
\noindent\textbf{Later retrieval:} Explain what a WPA2-Personal handshake demonstrates and what it does not hide.\par
\noindent\textbf{Alignment references:} htb-foundations, ies-practice\par
\subsubsection{networking-for-security: capstone}
\noindent\textbf{Milestone:} capstone-milestone\par
\noindent\textbf{Skill:} Diagnose an unfamiliar network fault by testing competing explanations.\par
\noindent\textbf{Experiment:} Use a fresh fault in a disposable lab or an unsolved training scenario. Record symptoms, at least two hypotheses, discriminating checks, root cause, and verified repair.\par
\noindent\textbf{Transfer:} Introduce or select a different fault with a similar symptom and distinguish it using evidence.\par
\noindent\textbf{Control:} Repeat both a formerly failing check and an unrelated healthy check after repair.\par
\noindent\textbf{Later retrieval:} From a new symptom, name the first two tests and what each possible result would rule out.\par
\noindent\textbf{Alignment references:} htb-foundations, ies-practice\par
\subsubsection{web-application-security: injection}
\noindent\textbf{Milestone:} injection-milestone\par
\noindent\textbf{Skill:} Infer an injection mechanism from request/response evidence and distinguish a real signal from noise.\par
\noindent\textbf{Experiment:} Solve an unfamiliar scoped Academy injection lab; save redacted baseline and successful requests, observations, root cause, and a specific repair.\par
\noindent\textbf{Transfer:} Solve a second small variation of the same mechanism with a different input location, data shape, or response signal.\par
\noindent\textbf{Control:} Include a negative request and explain why a status change or delay alone is insufficient evidence.\par
\noindent\textbf{Later retrieval:} Reconstruct the data-to-query or data-to-command boundary from memory and predict a new input.\par
\noindent\textbf{Alignment references:} portswigger-paths, portswigger-mystery, htb-foundations\par
\subsubsection{web-application-security: auth}
\noindent\textbf{Milestone:} auth-milestone\par
\noindent\textbf{Skill:} Track authentication state and identify the missing validation in an account-recovery or session flow.\par
\noindent\textbf{Experiment:} Solve a scoped multi-step authentication lab; record the state transition that incorrectly grants access and the intended authentication boundary.\par
\noindent\textbf{Transfer:} Change the account relationship or order of recovery/session operations in a comparable training case. Predict which state transition permits or denies access, then test without a walkthrough.\par
\noindent\textbf{Control:} Show a legitimate login and a denied unauthenticated request; specify which server-side check breaks the attack chain.\par
\noindent\textbf{Later retrieval:} Draw the relevant session or recovery state transitions without notes and identify the vulnerable transition.\par
\noindent\textbf{Alignment references:} portswigger-paths, portswigger-mystery, htb-foundations\par
\subsubsection{web-application-security: access-control}
\noindent\textbf{Milestone:} access-control-milestone\par
\noindent\textbf{Skill:} Test authorization separately from authentication.\par
\noindent\textbf{Experiment:} Using two lab users and an administrator role where available, build an owner/other-user/anonymous access matrix for one resource and demonstrate a missing check.\par
\noindent\textbf{Transfer:} Change the object identifier or request method and repeat the authorization test without following a walkthrough.\par
\noindent\textbf{Control:} Include the legitimate owner request and a correctly denied cross-user request or a repaired local endpoint.\par
\noindent\textbf{Later retrieval:} Explain why knowing an object ID does not grant authority, then design a fresh access test.\par
\noindent\textbf{Alignment references:} portswigger-paths, portswigger-mystery, htb-foundations\par
\subsubsection{web-application-security: xss-and-csrf}
\noindent\textbf{Milestone:} xss-csrf-milestone\par
\noindent\textbf{Skill:} Separate script execution caused by unsafe output handling from an unwanted authenticated cross-site action.\par
\noindent\textbf{Experiment:} Complete a scoped XSS exercise and a scoped CSRF exercise; trace source, browser context, and effect for each. Use harmless lab proofs.\par
\noindent\textbf{Transfer:} Change the XSS rendering context or the CSRF request conditions and predict whether the previous proof still works.\par
\noindent\textbf{Control:} Compare correctly encoded output or a rejected anti-CSRF token; explain why an XSS fix is not automatically a CSRF fix.\par
\noindent\textbf{Later retrieval:} From two request traces, distinguish XSS from CSRF and name the relevant trust boundary.\par
\noindent\textbf{Alignment references:} portswigger-paths, portswigger-mystery, htb-foundations\par
\subsubsection{web-application-security: ssrf-and-deserialization}
\noindent\textbf{Milestone:} ssrf-deser-milestone\par
\noindent\textbf{Skill:} Distinguish server-side fetching, external-entity expansion, and unsafe object reconstruction.\par
\noindent\textbf{Experiment:} Use scoped lab evidence from SSRF, XXE, and deserialization exercises to explain each input-to-action chain; independently reproduce one on an unfamiliar lab.\par
\noindent\textbf{Transfer:} Change the SSRF destination, XML entity setting, or serialized property and predict the result before testing in the lab.\par
\noindent\textbf{Control:} Include a rejected or inert input and identify the specific allowlist, parser, or data-format boundary that prevents the action.\par
\noindent\textbf{Later retrieval:} Classify three short scenarios by mechanism and explain why cloud metadata access is not implied by every SSRF.\par
\noindent\textbf{Alignment references:} portswigger-paths, portswigger-mystery, htb-foundations\par
\subsubsection{web-application-security: logic-flaws}
\noindent\textbf{Milestone:} logic-milestone\par
\noindent\textbf{Skill:} Express a business invariant and produce a sequence that violates it.\par
\noindent\textbf{Experiment:} Choose a lab workflow, write its expected invariant, and test reordered, repeated, or concurrent operations to find a violation.\par
\noindent\textbf{Transfer:} Change the operation order or concurrency and predict whether the invariant still fails.\par
\noindent\textbf{Control:} Show the ordinary sequence preserving the invariant and propose a server-side transaction or state check.\par
\noindent\textbf{Later retrieval:} State an invariant for a different workflow and propose a discriminating test.\par
\noindent\textbf{Alignment references:} portswigger-paths, portswigger-mystery, htb-foundations\par
\subsubsection{web-application-security: api-security}
\noindent\textbf{Milestone:} api-milestone\par
\noindent\textbf{Skill:} Test API authorization and token handling using a bounded endpoint inventory.\par
\noindent\textbf{Experiment:} Map a small lab API and test anonymous, owner, and other-user access for the selected operations. Explain one JWT or GraphQL finding with a reproducible request.\par
\noindent\textbf{Transfer:} Change an object ID, field, role, or token property and test whether the proposed boundary holds.\par
\noindent\textbf{Control:} Show a correctly rejected request; do not label introspection alone a vulnerability without sensitive exposure or unauthorized access.\par
\noindent\textbf{Later retrieval:} Design an owner/other-user API test from memory and explain the difference between token decoding and signature validation.\par
\noindent\textbf{Alignment references:} portswigger-paths, portswigger-mystery, htb-foundations\par
\subsubsection{web-application-security: capstone}
\noindent\textbf{Milestone:} capstone-milestone\par
\noindent\textbf{Skill:} Produce an assessment another learner can reproduce and critically review.\par
\noindent\textbf{Experiment:} On a fresh scoped target, reproduce three findings from distinct covered mechanisms. Report scope, evidence, impact, root cause, and a specific fix for each.\par
\noindent\textbf{Transfer:} Change a relevant request input, role, or workflow in an owned reproduction of one finding. Predict whether the same weakness applies and independently test the prediction.\par
\noindent\textbf{Control:} Implement and retest at least one fix in an owned local reproduction, keeping legitimate behavior working; label other fixes as proposed.\par
\noindent\textbf{Later retrieval:} Explain one finding and its fix from memory, then test a changed request.\par
\noindent\textbf{Alignment references:} portswigger-paths, portswigger-mystery, htb-foundations\par
\subsubsection{network-penetration-testing: recon}
\noindent\textbf{Milestone:} recon-milestone\par
\noindent\textbf{Skill:} Turn reconnaissance observations into falsifiable service hypotheses.\par
\noindent\textbf{Experiment:} Map a fresh owned/training host: ports, services, reachable web entries, scan limitations, and two ranked hypotheses before exploitation.\par
\noindent\textbf{Transfer:} Move a lab service to a nondefault port or change its banner. Infer its actual protocol from new observations and confirm it with a native client or second tool.\par
\noindent\textbf{Control:} Test a closed or filtered port and an intentionally nonexistent web route; distinguish them from positive findings.\par
\noindent\textbf{Later retrieval:} Explain the next check for a new banner, including what would falsify the suspected weakness.\par
\noindent\textbf{Alignment references:} htb-pentest, htb-foundations\par
\subsubsection{network-penetration-testing: service-exploitation}
\noindent\textbf{Milestone:} service-exploitation-milestone\par
\noindent\textbf{Skill:} Verify exploit prerequisites and explain the boundary crossed.\par
\noindent\textbf{Experiment:} On a designated vulnerable lab, establish a foothold using a documented service weakness after verifying its required conditions. Save sanitized evidence.\par
\noindent\textbf{Transfer:} Use a comparable training service with one changed configuration or prerequisite. Predict whether the technique still applies, adapt the test, and explain success or failure.\par
\noindent\textbf{Control:} Remove one prerequisite in an owned reproduction and show that the exploit fails while legitimate access remains.\par
\noindent\textbf{Later retrieval:} Explain why a matching CVE or version is a hypothesis, then enumerate the required conditions.\par
\noindent\textbf{Alignment references:} htb-pentest, htb-foundations\par
\subsubsection{network-penetration-testing: metasploit}
\noindent\textbf{Milestone:} metasploit-milestone\par
\noindent\textbf{Skill:} Separate framework orchestration from the underlying exploit mechanism.\par
\noindent\textbf{Experiment:} Complete one bounded training exploit with Metasploit, recording target assumptions, module options, payload role, and proof of the resulting context.\par
\noindent\textbf{Transfer:} Reproduce the essential request or underlying mechanism manually in the same lab, using earlier exercises as preparation.\par
\noindent\textbf{Control:} Use an incompatible target setting or a patched local reproduction and explain the failure rather than changing options blindly.\par
\noindent\textbf{Later retrieval:} Describe the exploit/payload distinction and reconstruct the required options without a saved session.\par
\noindent\textbf{Alignment references:} htb-pentest, htb-foundations\par
\subsubsection{network-penetration-testing: linux-privesc}
\noindent\textbf{Milestone:} linux-privesc-milestone\par
\noindent\textbf{Skill:} Find and validate Linux privilege boundaries through native permissions and process context.\par
\noindent\textbf{Experiment:} From a provided low-privilege lab foothold, demonstrate one escalation path and manually verify a second distinct candidate; use a second lab if necessary.\par
\noindent\textbf{Transfer:} Repeat the successful mechanism with a different lab file, executable, or scheduled task and predict the permission change.\par
\noindent\textbf{Control:} Correct the writable path, ownership, or privilege setting in a disposable reproduction and retest.\par
\noindent\textbf{Later retrieval:} Explain real/effective identity and one trust boundary from memory, then inspect a new permission example.\par
\noindent\textbf{Alignment references:} htb-pentest, htb-foundations\par
\subsubsection{network-penetration-testing: windows-privesc}
\noindent\textbf{Milestone:} windows-privesc-milestone\par
\noindent\textbf{Skill:} Reason about Windows identity, service permissions, and token privileges.\par
\noindent\textbf{Experiment:} Establish a baseline with native Windows commands, then complete one service or token-based lab escalation and explain the exact enabling permissions.\par
\noindent\textbf{Transfer:} Change the lab user or service configuration and predict whether the same path remains available.\par
\noindent\textbf{Control:} Show a denied operation or corrected service ACL; record actual identity and privileges rather than assuming every shell is SYSTEM.\par
\noindent\textbf{Later retrieval:} Explain the difference between group membership, an enabled privilege, and a service access-control entry.\par
\noindent\textbf{Alignment references:} htb-pentest, htb-foundations\par
\subsubsection{network-penetration-testing: active-directory}
\noindent\textbf{Milestone:} ad-milestone\par
\noindent\textbf{Skill:} Validate an AD attack-path edge using identity, permissions, and authentication evidence.\par
\noindent\textbf{Experiment:} In an isolated beginner AD lab with supplied credentials, enumerate the domain and validate a bounded path to the designated objective. Record each required edge.\par
\noindent\textbf{Transfer:} Re-evaluate one graph edge with another lab account or changed permission; verify the graph claim using native evidence.\par
\noindent\textbf{Control:} Remove or identify a specific permission or weak credential condition that breaks the path; do not claim all AD attacks require Kerberoasting.\par
\noindent\textbf{Later retrieval:} Draw a small domain trust/permission path from memory and explain a service ticket without equating it to a plaintext password.\par
\noindent\textbf{Alignment references:} htb-pentest, htb-foundations\par
\subsubsection{network-penetration-testing: post-exploitation}
\noindent\textbf{Milestone:} post-exploitation-milestone\par
\noindent\textbf{Skill:} Map observed actions to telemetry and distinguish visibility gaps from evidence of absence.\par
\noindent\textbf{Experiment:} Map the stages of a completed lab exercise to ATT\&CK techniques and actual log or network evidence, including a harmless scheduled-task persistence fixture and synthetic file transfer.\par
\noindent\textbf{Transfer:} Change the task name or transfer channel and predict how the evidence changes.\par
\noindent\textbf{Control:} Remove the lab artifact and confirm its absence; label an unlogged action as a visibility gap instead of claiming it did not happen.\par
\noindent\textbf{Later retrieval:} Explain which logs would confirm a new persistence or transfer hypothesis and what they cannot prove.\par
\noindent\textbf{Alignment references:} htb-pentest, htb-foundations\par
\subsubsection{network-penetration-testing: capstone}
\noindent\textbf{Milestone:} capstone-milestone\par
\noindent\textbf{Skill:} Complete and report a scoped assessment without following a target walkthrough.\par
\noindent\textbf{Experiment:} On an unfamiliar training box, progress from reconnaissance to the designated foothold and privilege objective, recording dead ends, evidence, and help used.\par
\noindent\textbf{Transfer:} In an owned reproduction, change one account permission, service setting, or route relevant to a key step. Predict the effect on the assessment chain and independently retest.\par
\noindent\textbf{Control:} Apply one hardening change in an owned reproduction and show it blocks that step without breaking intended use; clearly label untested fixes.\par
\noindent\textbf{Later retrieval:} Explain the assessment chain and one alternative hypothesis from memory, then rerun the discriminating check.\par
\noindent\textbf{Alignment references:} htb-pentest, htb-foundations\par
\subsubsection{reverse-engineering-and-binary-exploitation: asm-foundations}
\noindent\textbf{Milestone:} asm-milestone\par
\noindent\textbf{Skill:} Predict a function’s behavior from registers, stack, branches, and the target ABI.\par
\noindent\textbf{Experiment:} Trace a small compiled function for two inputs before running it, then verify the predicted return value and key state changes in a debugger.\par
\noindent\textbf{Transfer:} Change one argument or compile option and predict which instruction or state observation changes.\par
\noindent\textbf{Control:} Include an input taking the other branch; record ABI and architecture so register names are not treated as universal.\par
\noindent\textbf{Later retrieval:} Sketch the call/return state and argument locations for the chosen ABI without notes.\par
\noindent\textbf{Alignment references:} pwn-college, rop-guide, phoenix-stack0\par
\subsubsection{reverse-engineering-and-binary-exploitation: static-analysis}
\noindent\textbf{Milestone:} static-milestone\par
\noindent\textbf{Skill:} Recover a small program’s input condition from static evidence and verify it dynamically.\par
\noindent\textbf{Experiment:} Analyze an unfamiliar educational crackme without source; document the candidate input predicate and its assembly/decompiler evidence.\par
\noindent\textbf{Transfer:} Predict an accepted and a rejected input not both supplied by a walkthrough, then run them.\par
\noindent\textbf{Control:} Cross-check the decompiler against instructions; an unclear type or optimized-away function is uncertainty, not proof the decompiler is wrong.\par
\noindent\textbf{Later retrieval:} Explain the recovered predicate and locate its key comparison in a fresh small function.\par
\noindent\textbf{Alignment references:} pwn-college, rop-guide, phoenix-stack0\par
\subsubsection{reverse-engineering-and-binary-exploitation: dynamic-analysis}
\noindent\textbf{Milestone:} dynamic-milestone\par
\noindent\textbf{Skill:} Use controlled execution to explain an input-dependent state change.\par
\noindent\textbf{Experiment:} Solve a small educational crackme using breakpoints, register/memory inspection, and a trace of the deciding branch.\par
\noindent\textbf{Transfer:} Change the input and predict the next branch or watched value before continuing.\par
\noindent\textbf{Control:} Demonstrate a rejected input and distinguish your own debugger modifications from an input that works in an unmodified run.\par
\noindent\textbf{Later retrieval:} Choose a breakpoint/watchpoint for a new toy function and predict what it will reveal.\par
\noindent\textbf{Alignment references:} pwn-college, rop-guide, phoenix-stack0\par
\subsubsection{reverse-engineering-and-binary-exploitation: stack-overflows}
\noindent\textbf{Milestone:} stack-milestone\par
\noindent\textbf{Skill:} Explain how an out-of-bounds write progresses from data corruption to control-flow corruption.\par
\noindent\textbf{Experiment:} Use the Phoenix progression or an equivalent deliberately vulnerable local binary to demonstrate adjacent-data corruption and a later controlled redirection. Explain offset, ABI, and each payload region.\par
\noindent\textbf{Transfer:} Change padding or a small source buffer size in your own toy reproduction and recalculate the relevant boundary.\par
\noindent\textbf{Control:} Run a bounded-input version or a short input and show that the control target is not overwritten.\par
\noindent\textbf{Later retrieval:} Draw the relevant stack layout and distinguish Phoenix Stack Zero’s data overwrite from overwriting a saved return address.\par
\noindent\textbf{Alignment references:} pwn-college, rop-guide, phoenix-stack0\par
\subsubsection{reverse-engineering-and-binary-exploitation: protections}
\noindent\textbf{Milestone:} protections-milestone\par
\noindent\textbf{Skill:} Distinguish mitigation mechanisms instead of treating checksec as a proof of safety.\par
\noindent\textbf{Experiment:} Compile one owned toy program in four protection configurations and compare flags with actual addresses or failure behavior.\par
\noindent\textbf{Transfer:} Change one protection at a time and predict which earlier technique is affected.\par
\noindent\textbf{Control:} Include a bounds-corrected build: mitigations and fixing the memory bug are different interventions.\par
\noindent\textbf{Later retrieval:} Explain ASLR, PIE, canaries, and NX from memory and name a limitation of each.\par
\noindent\textbf{Alignment references:} pwn-college, rop-guide, phoenix-stack0\par
\subsubsection{reverse-engineering-and-binary-exploitation: rop}
\noindent\textbf{Milestone:} rop-milestone\par
\noindent\textbf{Skill:} Construct and explain a chain of existing instructions under a specific calling convention.\par
\noindent\textbf{Experiment:} Complete the ROP Emporium callme challenge, or an equivalent challenge requiring several function calls; annotate arguments, gadgets, stack effects, and alignment.\par
\noindent\textbf{Transfer:} Change one argument or gadget choice and predict whether the calls still satisfy the challenge.\par
\noindent\textbf{Control:} Show one incorrect chain and explain its precise failure; ret2win alone is not evidence of multi-call chaining.\par
\noindent\textbf{Later retrieval:} Reconstruct a short gadget chain and register setup from memory for the chosen architecture.\par
\noindent\textbf{Alignment references:} pwn-college, rop-guide, phoenix-stack0\par
\subsubsection{reverse-engineering-and-binary-exploitation: heap-and-format}
\noindent\textbf{Milestone:} heap-format-milestone\par
\noindent\textbf{Skill:} Distinguish allocation lifetime errors from format-string interpretation bugs.\par
\noindent\textbf{Experiment:} Use prior local exercises to show a heap lifetime error and a format-string information leak, then independently solve one introductory lab in either class.\par
\noindent\textbf{Transfer:} Change an allocation/use order or format input and predict the observable result.\par
\noindent\textbf{Control:} Compare a lifetime-correct or fixed-format version; allocator behavior depends on version and undefined behavior need not crash.\par
\noindent\textbf{Later retrieval:} Classify a new short example and explain the relevant object lifetime or variadic format boundary.\par
\noindent\textbf{Alignment references:} pwn-college, rop-guide, phoenix-stack0\par
\subsubsection{reverse-engineering-and-binary-exploitation: capstone}
\noindent\textbf{Milestone:} capstone-milestone\par
\noindent\textbf{Skill:} Move from a defensible vulnerability hypothesis to a reproducible lab exploit and repair.\par
\noindent\textbf{Experiment:} Choose a fresh educational binary matched to your architecture and level. Record static/dynamic evidence, vulnerability mechanism, exploit, and the actual enabled protections.\par
\noindent\textbf{Transfer:} Change one relevant input, source buffer size, or mitigation in an owned equivalent reproduction. Predict whether the exploit still applies, then test and explain the changed assumptions.\par
\noindent\textbf{Control:} Patch an equivalent owned source reproduction and verify both rejection of the bad input and preservation of intended behavior.\par
\noindent\textbf{Later retrieval:} Explain the vulnerability-to-effect chain without the exploit script, then predict a changed input.\par
\noindent\textbf{Alignment references:} pwn-college, rop-guide, phoenix-stack0\par
\subsubsection{defensive-security-and-detection: logging-fundamentals}
\noindent\textbf{Milestone:} logging-milestone\par
\noindent\textbf{Skill:} Reconstruct an event timeline while separating observations from inference.\par
\noindent\textbf{Experiment:} Use three benign events from your lab or an unused training log. Reconstruct order, identity, action, and outcome from source lines before comparing with ground truth.\par
\noindent\textbf{Transfer:} Change the time window or include another source and explain timestamp/timezone handling.\par
\noindent\textbf{Control:} Include an unrelated event and one deliberately missing source; distinguish a negative result from missing telemetry.\par
\noindent\textbf{Later retrieval:} From memory, explain which fields connect an authentication event to a host and time.\par
\noindent\textbf{Alignment references:} htb-soc, security-onion-hardware, ies-practice\par
\subsubsection{defensive-security-and-detection: log-analysis}
\noindent\textbf{Milestone:} log-analysis-milestone\par
\noindent\textbf{Skill:} Build a detection query and evaluate it against labeled positive and negative data.\par
\noindent\textbf{Experiment:} Use a small labeled training dataset containing suspicious and ordinary requests. Design a query without looking at the answer labels, then compare its results with ground truth.\par
\noindent\textbf{Transfer:} Change the source address or request ordering in a synthetic copy; test whether the same behavior is still found.\par
\noindent\textbf{Control:} Record true positives, false positives, and false negatives with counts and explain at least one ambiguity.\par
\noindent\textbf{Later retrieval:} Write the core filter from memory and test it on a small changed dataset.\par
\noindent\textbf{Alignment references:} htb-soc, security-onion-hardware, ies-practice\par
\subsubsection{defensive-security-and-detection: network-monitoring}
\noindent\textbf{Milestone:} network-monitoring-milestone\par
\noindent\textbf{Skill:} Explain the differences between connection metadata and signature alerts.\par
\noindent\textbf{Experiment:} Replay a documented benign-plus-test PCAP through Zeek and Suricata; confirm the relevant parser and rule are enabled before comparing outputs.\par
\noindent\textbf{Transfer:} Change a protocol field or replay a different small capture and predict the connection record and alert outcome.\par
\noindent\textbf{Control:} Use a benign control capture; a generic nmap scan is not guaranteed to trigger an arbitrary default ruleset.\par
\noindent\textbf{Later retrieval:} Explain why a connection log can exist without an alert, then identify the required visibility and rule.\par
\noindent\textbf{Alignment references:} htb-soc, security-onion-hardware, ies-practice\par
\subsubsection{defensive-security-and-detection: siem-basics}
\noindent\textbf{Milestone:} siem-milestone\par
\noindent\textbf{Skill:} Investigate indexed security events using reproducible queries and pivots.\par
\noindent\textbf{Experiment:} In Security Onion Import/Eval on supported x86-64 hardware, or a supplied hosted SIEM lab, ingest/inspect a small training dataset and investigate an unfamiliar event.\par
\noindent\textbf{Transfer:} Pivot from an event to a second host, identity, or time interval without using a prewritten query for that pivot.\par
\noindent\textbf{Control:} Confirm a known benign event is also ingested; account for missing sources, time offsets, and query scope before inferring absence.\par
\noindent\textbf{Later retrieval:} Rebuild the investigation query and explain its fields without using the saved dashboard.\par
\noindent\textbf{Alignment references:} htb-soc, security-onion-hardware, ies-practice\par
\subsubsection{defensive-security-and-detection: detection-engineering}
\noindent\textbf{Milestone:} detection-milestone\par
\noindent\textbf{Skill:} Evaluate a rule as a hypothesis about observable behavior.\par
\noindent\textbf{Experiment:} Write one new rule for a bounded lab behavior; record required telemetry, test data, rule version, and results against labeled positive and negative samples.\par
\noindent\textbf{Transfer:} Change a nonessential address, username, or ordering and check whether behavior-based coverage survives.\par
\noindent\textbf{Control:} Count false positives and false negatives; explain a known blind spot and test the benign control after tuning.\par
\noindent\textbf{Later retrieval:} Explain the rule’s observable assumption from memory and design one counterexample.\par
\noindent\textbf{Alignment references:} htb-soc, security-onion-hardware, ies-practice\par
\subsubsection{defensive-security-and-detection: hardening}
\noindent\textbf{Milestone:} hardening-milestone\par
\noindent\textbf{Skill:} Demonstrate that a configuration change blocks a chosen path while preserving legitimate use.\par
\noindent\textbf{Experiment:} On a disposable VM, establish baseline access, apply selected service/permission/update-policy changes, and rerun the relevant lab checks.\par
\noindent\textbf{Transfer:} Test from a second role or service endpoint to find an overbroad assumption.\par
\noindent\textbf{Control:} Prove legitimate access still works and verify rollback. An empty scan alone does not distinguish hardening from an offline machine.\par
\noindent\textbf{Later retrieval:} Explain one hardening control, its operational tradeoff, and a positive/negative test pair.\par
\noindent\textbf{Alignment references:} htb-soc, security-onion-hardware, ies-practice\par
\subsubsection{defensive-security-and-detection: incident-response}
\noindent\textbf{Milestone:} ir-milestone\par
\noindent\textbf{Skill:} Make evidence-preserving response decisions under uncertainty.\par
\noindent\textbf{Experiment:} Use an unfamiliar training case or a peer-run benign lab scenario. Document timeline, evidence hashes, containment choices, investigation, and recovery checks before opening the answer key.\par
\noindent\textbf{Transfer:} Consider one new piece of evidence that contradicts the initial hypothesis and revise the response.\par
\noindent\textbf{Control:} Show that collection copies remain unchanged and that recovery restores intended service; label simulated containment actions as simulated.\par
\noindent\textbf{Later retrieval:} Explain the first response steps for a new symptom and what could be lost through an immediate reboot.\par
\noindent\textbf{Alignment references:} htb-soc, security-onion-hardware, ies-practice\par
\subsubsection{defensive-security-and-detection: capstone}
\noindent\textbf{Milestone:} capstone-milestone\par
\noindent\textbf{Skill:} Measure monitoring coverage and close a bounded detection gap without claiming complete visibility.\par
\noindent\textbf{Experiment:} Observe or replay a small documented lab chain, compare expected and observed telemetry, and deliver a reproducible investigation with coverage and gaps.\par
\noindent\textbf{Transfer:} Replay a varied form of one stage and evaluate whether the improved rule still detects it.\par
\noindent\textbf{Control:} Retest benign traffic after tuning and record residual gaps that require telemetry changes rather than another rule.\par
\noindent\textbf{Later retrieval:} Reconstruct the incident and one gap from memory, then explain how a new analyst could reproduce the key query.\par
\noindent\textbf{Alignment references:} htb-soc, security-onion-hardware, ies-practice\par
\subsubsection{ethical-hacking-foundations: foundations}
\noindent\textbf{Milestone:} foundations-milestone\par
\noindent\textbf{Skill:} Prove the entry skills in a small lab\par
\noindent\textbf{Experiment:} Demonstrate each entry skill with a small owned fixture or permitted training environment. Record the actual setup, one observation per skill and any prerequisite you need to repair.\par
\noindent\textbf{Transfer:} Change a relevant identity, input or environment setting and predict the effect before repeating the demonstration.\par
\noindent\textbf{Control:} Include a deliberately unmet prerequisite and distinguish setup failure from a result about the subject.\par
\noindent\textbf{Later retrieval:} Reconstruct the environment and entry assumptions without following the setup notes.\par
\noindent\textbf{Alignment references:} offsec-pen, htb-foundations, portswigger-paths\par
\subsubsection{ethical-hacking-foundations: scope}
\noindent\textbf{Milestone:} scope-milestone\par
\noindent\textbf{Skill:} What authorizes this test?\par
\noindent\textbf{Experiment:} Write rules of engagement for a loopback training app, including allowed actions, rate limits, stop conditions, contacts and evidence handling.\par
\noindent\textbf{Transfer:} Add a third-party login dependency and decide which requests must remain excluded until permission changes.\par
\noindent\textbf{Control:} Show a technically reachable endpoint that the written scope forbids, and leave it untouched.\par
\noindent\textbf{Later retrieval:} Reconstruct the difference between technical access, ownership and authorization.\par
\noindent\textbf{Alignment references:} offsec-pen, htb-foundations, portswigger-paths\par
\subsubsection{ethical-hacking-foundations: method}
\noindent\textbf{Milestone:} method-milestone\par
\noindent\textbf{Skill:} Follow a hypothesis, not a tool list\par
\noindent\textbf{Experiment:} Inventory the disposable app, rank two vulnerability hypotheses, and run a small non-destructive check for each with sanitized request/response evidence.\par
\noindent\textbf{Transfer:} Move the app to a new port or change a visible header and revise the hypotheses using actual behavior.\par
\noindent\textbf{Control:} Compare a real route with an intentionally nonexistent route so a generic error page is not counted as a finding.\par
\noindent\textbf{Later retrieval:} Given a new banner, name the next observation that could disprove your hypothesis.\par
\noindent\textbf{Alignment references:} offsec-pen, htb-foundations, portswigger-paths\par
\subsubsection{ethical-hacking-foundations: web-boundaries}
\noindent\textbf{Milestone:} web-boundaries-milestone\par
\noindent\textbf{Skill:} Two accounts, one boundary\par
\noindent\textbf{Experiment:} In an explicitly designated authorization lesson, test an owner and a second account against the same synthetic object and explain the server-side boundary.\par
\noindent\textbf{Transfer:} Change the object or HTTP method and predict whether the authorization decision should change.\par
\noindent\textbf{Control:} Demonstrate legitimate owner access and a denied unauthorized request on a corrected local equivalent.\par
\noindent\textbf{Later retrieval:} Explain why hiding a button is different from enforcing authorization.\par
\noindent\textbf{Alignment references:} offsec-pen, htb-foundations, portswigger-paths\par
\subsubsection{ethical-hacking-foundations: permissions}
\noindent\textbf{Milestone:} permissions-milestone\par
\noindent\textbf{Skill:} Who is this process allowed to be?\par
\noindent\textbf{Experiment:} Create two local test roles or use a supplied permissions fixture, identify an overbroad grant, and repair it while recording before/after effective access.\par
\noindent\textbf{Transfer:} Change group membership or parent-directory permissions and predict the new access result.\par
\noindent\textbf{Control:} Keep the intended reader working and demonstrate that the unrelated role remains denied.\par
\noindent\textbf{Later retrieval:} Explain the difference between authentication, authorization and process privilege.\par
\noindent\textbf{Alignment references:} offsec-pen, htb-foundations, portswigger-paths\par
\subsubsection{ethical-hacking-foundations: evidence}
\noindent\textbf{Milestone:} evidence-milestone\par
\noindent\textbf{Skill:} Make a finding another person can verify\par
\noindent\textbf{Experiment:} Write one lab finding with exact target/version, minimal steps, sanitized evidence, severity rationale and a tested local repair.\par
\noindent\textbf{Transfer:} Give a peer a changed input or repeat on a second clean fixture and identify assumptions missing from the report.\par
\noindent\textbf{Control:} Show that the fix rejects the bad case without breaking an ordinary valid request.\par
\noindent\textbf{Later retrieval:} Reconstruct a finding from its root cause without copying the original prose.\par
\noindent\textbf{Alignment references:} offsec-pen, htb-foundations, portswigger-paths\par
\subsubsection{ethical-hacking-foundations: assessment}
\noindent\textbf{Milestone:} assessment-milestone\par
\noindent\textbf{Skill:} Run a small authorized assessment\par
\noindent\textbf{Experiment:} Perform scope confirmation, inventory, two hypothesis tests and one reproducible finding or justified negative result; deliver an executive summary and technical evidence.\par
\noindent\textbf{Transfer:} Change an account role or route and independently reassess whether the suspected weakness remains.\par
\noindent\textbf{Control:} Include an out-of-scope action you deliberately excluded and a benign workflow that survives the repair.\par
\noindent\textbf{Later retrieval:} Explain the assessment chain, one false lead and what evidence changed your mind.\par
\noindent\textbf{Alignment references:} offsec-pen, htb-foundations, portswigger-paths\par
\subsubsection{ethical-hacking-foundations: independent-transfer}
\noindent\textbf{Milestone:} independent-transfer-milestone\par
\noindent\textbf{Skill:} Revisit the skill on a fresh problem\par
\noindent\textbf{Experiment:} After a delay, attempt the fresh task and produce a reproducible evidence packet. Ask a peer to select or inspect the variation when possible; record assistance, failed hypotheses and remaining limits.\par
\noindent\textbf{Transfer:} Change a mechanism-relevant condition, rather than only resetting the old environment, and independently predict and test the resulting behavior.\par
\noindent\textbf{Control:} Include a contrasting case where the mechanism should not apply and a legitimate operation that should remain valid.\par
\noindent\textbf{Later retrieval:} Explain what transferred, what required help and which prerequisite needs another small experiment.\par
\noindent\textbf{Alignment references:} offsec-pen, htb-foundations, portswigger-paths\par
\subsubsection{red-team-adversary-emulation: foundations}
\noindent\textbf{Milestone:} foundations-milestone\par
\noindent\textbf{Skill:} Prove the entry skills in a small lab\par
\noindent\textbf{Experiment:} Demonstrate each entry skill with a small owned fixture or permitted training environment. Record the actual setup, one observation per skill and any prerequisite you need to repair.\par
\noindent\textbf{Transfer:} Change a relevant identity, input or environment setting and predict the effect before repeating the demonstration.\par
\noindent\textbf{Control:} Include a deliberately unmet prerequisite and distinguish setup failure from a result about the subject.\par
\noindent\textbf{Later retrieval:} Reconstruct the environment and entry assumptions without following the setup notes.\par
\noindent\textbf{Alignment references:} mitre-emulation, atomic-red, htb-pentest\par
\subsubsection{red-team-adversary-emulation: objectives}
\noindent\textbf{Milestone:} objectives-milestone\par
\noindent\textbf{Skill:} Measure an objective, not a shell\par
\noindent\textbf{Experiment:} Author an engagement plan with rules of engagement, white-cell contacts, allowed assets, deconfliction, stop conditions and a defender measurement plan.\par
\noindent\textbf{Transfer:} Add an operational constraint such as unavailable endpoint logging and redesign the objective or telemetry plan honestly.\par
\noindent\textbf{Control:} Show a scenario that appears successful technically but fails the stated business objective.\par
\noindent\textbf{Later retrieval:} State the difference between a penetration-test finding and an adversary-emulation objective.\par
\noindent\textbf{Alignment references:} mitre-emulation, atomic-red, htb-pentest\par
\subsubsection{red-team-adversary-emulation: threat-model}
\noindent\textbf{Milestone:} threat-model-milestone\par
\noindent\textbf{Skill:} Choose techniques from evidence\par
\noindent\textbf{Experiment:} Build a small ATT\&CK-aligned hypothesis from a cited threat report, specifying prerequisites, observable effects and a harmless substitute for each risky action.\par
\noindent\textbf{Transfer:} Replace the assumed identity or endpoint type and remove techniques whose prerequisites no longer hold.\par
\noindent\textbf{Control:} Include one plausible but unsupported technique and explain why it was excluded.\par
\noindent\textbf{Later retrieval:} Rebuild the evidence-to-technique chain and name its uncertainty.\par
\noindent\textbf{Alignment references:} mitre-emulation, atomic-red, htb-pentest\par
\subsubsection{red-team-adversary-emulation: telemetry}
\noindent\textbf{Milestone:} telemetry-milestone\par
\noindent\textbf{Skill:} Know what the defender can see\par
\noindent\textbf{Experiment:} Build a timeline joining the marker, account, process and network observations; label each missing source rather than inventing coverage.\par
\noindent\textbf{Transfer:} Change the marker path or process argument and predict which detection features remain stable.\par
\noindent\textbf{Control:} Repeat an ordinary administration action and distinguish its evidence from the staged behavior.\par
\noindent\textbf{Later retrieval:} Explain which sensor would be needed to resolve a missing observation.\par
\noindent\textbf{Alignment references:} mitre-emulation, atomic-red, htb-pentest\par
\subsubsection{red-team-adversary-emulation: emulation}
\noindent\textbf{Milestone:} emulation-milestone\par
\noindent\textbf{Skill:} Make a small emulation reversible\par
\noindent\textbf{Experiment:} In a disposable authorized lab, run a reviewed marker-only micro-emulation, record setup and exact effects, then perform and verify cleanup.\par
\noindent\textbf{Transfer:} Change a nonessential filename or account context and predict whether the intended behavior is still observable.\par
\noindent\textbf{Control:} Verify the artifact did not exist before, appears during the test, and is absent after cleanup; preserve unrelated files.\par
\noindent\textbf{Later retrieval:} Reconstruct the prerequisites and rollback steps without running a script blindly.\par
\noindent\textbf{Alignment references:} mitre-emulation, atomic-red, htb-pentest\par
\subsubsection{red-team-adversary-emulation: purple-team}
\noindent\textbf{Milestone:} purple-team-milestone\par
\noindent\textbf{Skill:} Test detection with a partner\par
\noindent\textbf{Experiment:} Have a defender or second review pass classify shuffled benign and emulated events; record misses, false positives and one revised analytic.\par
\noindent\textbf{Transfer:} Change a superficial indicator while preserving behavior and retest the revised analytic.\par
\noindent\textbf{Control:} Include an administration lookalike that should remain unalerted or receive a documented lower-confidence classification.\par
\noindent\textbf{Later retrieval:} Explain why a detection that matches only the demonstration filename is fragile.\par
\noindent\textbf{Alignment references:} mitre-emulation, atomic-red, htb-pentest\par
\subsubsection{red-team-adversary-emulation: campaign}
\noindent\textbf{Milestone:} campaign-milestone\par
\noindent\textbf{Skill:} Deliver a defensible exercise report\par
\noindent\textbf{Experiment:} Run the scoped sequence in a resettable lab, measure defender visibility and response, and deliver an objective-based report with technical evidence and remediation retests.\par
\noindent\textbf{Transfer:} Withhold one assumed permission or telemetry source and adapt the plan without exceeding scope.\par
\noindent\textbf{Control:} Demonstrate reset to baseline and distinguish untested coverage from a demonstrated miss.\par
\noindent\textbf{Later retrieval:} Explain what the exercise supports about defensive readiness and what it cannot establish.\par
\noindent\textbf{Alignment references:} mitre-emulation, atomic-red, htb-pentest\par
\subsubsection{red-team-adversary-emulation: independent-transfer}
\noindent\textbf{Milestone:} independent-transfer-milestone\par
\noindent\textbf{Skill:} Revisit the skill on a fresh problem\par
\noindent\textbf{Experiment:} After a delay, attempt the fresh task and produce a reproducible evidence packet. Ask a peer to select or inspect the variation when possible; record assistance, failed hypotheses and remaining limits.\par
\noindent\textbf{Transfer:} Change a mechanism-relevant condition, rather than only resetting the old environment, and independently predict and test the resulting behavior.\par
\noindent\textbf{Control:} Include a contrasting case where the mechanism should not apply and a legitimate operation that should remain valid.\par
\noindent\textbf{Later retrieval:} Explain what transferred, what required help and which prerequisite needs another small experiment.\par
\noindent\textbf{Alignment references:} mitre-emulation, atomic-red, htb-pentest\par
\subsubsection{soc-analyst: foundations}
\noindent\textbf{Milestone:} foundations-milestone\par
\noindent\textbf{Skill:} Prove the entry skills in a small lab\par
\noindent\textbf{Experiment:} Demonstrate each entry skill with a small owned fixture or permitted training environment. Record the actual setup, one observation per skill and any prerequisite you need to repair.\par
\noindent\textbf{Transfer:} Change a relevant identity, input or environment setting and predict the effect before repeating the demonstration.\par
\noindent\textbf{Control:} Include a deliberately unmet prerequisite and distinguish setup failure from a result about the subject.\par
\noindent\textbf{Later retrieval:} Reconstruct the environment and entry assumptions without following the setup notes.\par
\noindent\textbf{Alignment references:} microsoft-soc, htb-soc, cisa-kev\par
\subsubsection{soc-analyst: intake}
\noindent\textbf{Milestone:} intake-milestone\par
\noindent\textbf{Skill:} Turn an alert into a question\par
\noindent\textbf{Experiment:} Triage a synthetic queue containing a login anomaly, a web alert and a harmless maintenance event; assign urgency using evidence and business context.\par
\noindent\textbf{Transfer:} Change an asset from a sandbox to a critical service and justify whether prioritization changes.\par
\noindent\textbf{Control:} Keep a noisy but low-impact event from outranking a corroborated compromise solely because of event volume.\par
\noindent\textbf{Later retrieval:} Explain event, alert, incident, severity and confidence using one example.\par
\noindent\textbf{Alignment references:} microsoft-soc, htb-soc, cisa-kev\par
\subsubsection{soc-analyst: queries}
\noindent\textbf{Milestone:} queries-milestone\par
\noindent\textbf{Skill:} Follow the entity across logs\par
\noindent\textbf{Experiment:} Build reproducible queries over a supplied dataset to trace an identity, host and source address through a bounded time window, recording missing fields.\par
\noindent\textbf{Transfer:} Change the time zone or host identifier in a synthetic copy and adjust the query without losing the same activity.\par
\noindent\textbf{Control:} Include a similarly named account and an event outside the window that must not be joined.\par
\noindent\textbf{Later retrieval:} Recreate the query logic from the investigation question rather than memorizing syntax.\par
\noindent\textbf{Alignment references:} microsoft-soc, htb-soc, cisa-kev\par
\subsubsection{soc-analyst: authentication}
\noindent\textbf{Milestone:} authentication-milestone\par
\noindent\textbf{Skill:} Investigate a suspicious sign-in\par
\noindent\textbf{Experiment:} Investigate a provided authentication timeline, rank account misuse against benign alternatives, and identify the next discriminating evidence.\par
\noindent\textbf{Transfer:} Add a service-account explanation or a contradictory device event and update the conclusion.\par
\noindent\textbf{Control:} Include known maintenance logins and avoid treating location or failure count alone as compromise.\par
\noindent\textbf{Later retrieval:} Explain which evidence would raise or lower confidence in account misuse.\par
\noindent\textbf{Alignment references:} microsoft-soc, htb-soc, cisa-kev\par
\subsubsection{soc-analyst: web-and-mail}
\noindent\textbf{Milestone:} web-and-mail-milestone\par
\noindent\textbf{Skill:} Preserve evidence before following a lure\par
\noindent\textbf{Experiment:} Correlate an inert phishing-style message fixture with proxy/DNS and application logs, recording user exposure separately from demonstrated execution.\par
\noindent\textbf{Transfer:} Change a redirect or add a blocked request and reassess how far the activity progressed.\par
\noindent\textbf{Control:} Include a legitimate message with similar wording and a blocked attempt with no evidence of execution.\par
\noindent\textbf{Later retrieval:} Explain why receipt, click, execution and compromise are separate claims.\par
\noindent\textbf{Alignment references:} microsoft-soc, htb-soc, cisa-kev\par
\subsubsection{soc-analyst: response}
\noindent\textbf{Milestone:} response-milestone\par
\noindent\textbf{Skill:} Make containment proportionate\par
\noindent\textbf{Experiment:} Write and rehearse a lab incident runbook that preserves evidence, escalates uncertainty, contains a synthetic account or endpoint, and verifies intended access afterward.\par
\noindent\textbf{Transfer:} Make the affected system unavailable to the responder and choose an alternative action with explicit residual risk.\par
\noindent\textbf{Control:} Show that an unrelated account or service still works after containment.\par
\noindent\textbf{Later retrieval:} State the evidence needed before choosing a disruptive action.\par
\noindent\textbf{Alignment references:} microsoft-soc, htb-soc, cisa-kev\par
\subsubsection{soc-analyst: shift-handoff}
\noindent\textbf{Milestone:} shift-handoff-milestone\par
\noindent\textbf{Skill:} Close the case so someone can reopen it\par
\noindent\textbf{Experiment:} Investigate an unfamiliar small log bundle and deliver a timeline, evidence references, confidence, response decision and shift handoff with unresolved questions.\par
\noindent\textbf{Transfer:} Add a late contradictory event and revise the disposition while retaining the earlier reasoning.\par
\noindent\textbf{Control:} Document one rejected hypothesis and a known-benign comparison; do not turn missing logs into a clean bill of health.\par
\noindent\textbf{Later retrieval:} Reconstruct the case and next action from the handoff alone after a delay.\par
\noindent\textbf{Alignment references:} microsoft-soc, htb-soc, cisa-kev\par
\subsubsection{soc-analyst: independent-transfer}
\noindent\textbf{Milestone:} independent-transfer-milestone\par
\noindent\textbf{Skill:} Revisit the skill on a fresh problem\par
\noindent\textbf{Experiment:} After a delay, attempt the fresh task and produce a reproducible evidence packet. Ask a peer to select or inspect the variation when possible; record assistance, failed hypotheses and remaining limits.\par
\noindent\textbf{Transfer:} Change a mechanism-relevant condition, rather than only resetting the old environment, and independently predict and test the resulting behavior.\par
\noindent\textbf{Control:} Include a contrasting case where the mechanism should not apply and a legitimate operation that should remain valid.\par
\noindent\textbf{Later retrieval:} Explain what transferred, what required help and which prerequisite needs another small experiment.\par
\noindent\textbf{Alignment references:} microsoft-soc, htb-soc, cisa-kev\par
\subsubsection{active-directory-security: foundations}
\noindent\textbf{Milestone:} foundations-milestone\par
\noindent\textbf{Skill:} Prove the entry skills in a small lab\par
\noindent\textbf{Experiment:} Demonstrate each entry skill with a small owned fixture or permitted training environment. Record the actual setup, one observation per skill and any prerequisite you need to repair.\par
\noindent\textbf{Transfer:} Change a relevant identity, input or environment setting and predict the effect before repeating the demonstration.\par
\noindent\textbf{Control:} Include a deliberately unmet prerequisite and distinguish setup failure from a result about the subject.\par
\noindent\textbf{Later retrieval:} Reconstruct the environment and entry assumptions without following the setup notes.\par
\noindent\textbf{Alignment references:} microsoft-ad, microsoft-adcs, goad, htb-pentest\par
\subsubsection{active-directory-security: domain-model}
\noindent\textbf{Milestone:} domain-model-milestone\par
\noindent\textbf{Skill:} Draw the domain before testing it\par
\noindent\textbf{Experiment:} On a training domain, document DNS discovery, one domain join and a normal domain sign-in; preserve an architecture diagram and baseline event evidence.\par
\noindent\textbf{Transfer:} Change the client DNS server in the disposable lab and predict which discovery or authentication step fails.\par
\noindent\textbf{Control:} Compare a local login with a domain login and explain which identity authority answered.\par
\noindent\textbf{Later retrieval:} Recreate the domain discovery and trust boundaries from memory.\par
\noindent\textbf{Alignment references:} microsoft-ad, microsoft-adcs, goad, htb-pentest\par
\subsubsection{active-directory-security: auth-protocols}
\noindent\textbf{Milestone:} auth-protocols-milestone\par
\noindent\textbf{Skill:} A ticket is not a password\par
\noindent\textbf{Experiment:} Capture permitted authentication telemetry in a training domain and correlate ticket issuance or fallback with a controlled client action, without collecting real credentials.\par
\noindent\textbf{Transfer:} Change a service name or introduce documented clock skew in the disposable lab and predict the failure or fallback.\par
\noindent\textbf{Control:} Include a normal successful service access and a deliberately invalid service reference.\par
\noindent\textbf{Later retrieval:} Explain why DNS, SPNs and time synchronization matter to Kerberos.\par
\noindent\textbf{Alignment references:} microsoft-ad, microsoft-adcs, goad, htb-pentest\par
\subsubsection{active-directory-security: permissions}
\noindent\textbf{Milestone:} permissions-milestone\par
\noindent\textbf{Skill:} Prove an identity graph edge\par
\noindent\textbf{Experiment:} Trace a least-privilege violation in a supplied or owned directory graph and verify the claimed edge against native permissions before proposing a fix.\par
\noindent\textbf{Transfer:} Remove one group membership or inherited permission and predict whether the path still exists.\par
\noindent\textbf{Control:} Include a graph edge that appears plausible but is denied by the effective permissions.\par
\noindent\textbf{Later retrieval:} Explain the distinction between a graph hypothesis and a verified authorization path.\par
\noindent\textbf{Alignment references:} microsoft-ad, microsoft-adcs, goad, htb-pentest\par
\subsubsection{active-directory-security: gpo-delegation}
\noindent\textbf{Milestone:} gpo-delegation-milestone\par
\noindent\textbf{Skill:} Separate management power from service identity\par
\noindent\textbf{Experiment:} Inspect a training GPO and a documented delegation configuration, explain their intended function and identify an overbroad grant using effective access and normal service behavior.\par
\noindent\textbf{Transfer:} Restrict the policy link or delegated target and predict which administrative or service action stops working.\par
\noindent\textbf{Control:} Keep an intended management action functioning after the proposed least-privilege change.\par
\noindent\textbf{Later retrieval:} Explain why controlling a policy object and applying that policy are separate conditions.\par
\noindent\textbf{Alignment references:} microsoft-ad, microsoft-adcs, goad, htb-pentest\par
\subsubsection{active-directory-security: certificate-services}
\noindent\textbf{Milestone:} certificate-services-milestone\par
\noindent\textbf{Skill:} Follow the certificate enrollment decision\par
\noindent\textbf{Experiment:} In an isolated AD CS lab or supplied configuration exercise, map template enrollment permissions and identity settings to a concrete authentication risk; test a corrected configuration in a real lab before claiming operational completion.\par
\noindent\textbf{Transfer:} Change enrollment rights or an identity-related template setting and predict which request is permitted or rejected.\par
\noindent\textbf{Control:} Show a legitimate enrollment still works and an unauthorized identity request is rejected after the repair.\par
\noindent\textbf{Later retrieval:} Explain why possessing a valid certificate is not sufficient proof of an authorized identity.\par
\noindent\textbf{Alignment references:} microsoft-ad, microsoft-adcs, goad, htb-pentest\par
\subsubsection{active-directory-security: assessment}
\noindent\textbf{Milestone:} assessment-milestone\par
\noindent\textbf{Skill:} Harden a small domain and retest it\par
\noindent\textbf{Experiment:} Audit the training domain's discovery, authentication, permissions and one certificate or policy boundary; deliver native evidence, impact and a verified least-privilege repair.\par
\noindent\textbf{Transfer:} Change one account's memberships or a policy scope and independently re-evaluate the affected path.\par
\noindent\textbf{Control:} Retest normal user/service sign-in and document which enterprise controls were not present in the lab.\par
\noindent\textbf{Later retrieval:} Reconstruct one privilege path and the exact permission change that removed it.\par
\noindent\textbf{Alignment references:} microsoft-ad, microsoft-adcs, goad, htb-pentest\par
\subsubsection{active-directory-security: independent-transfer}
\noindent\textbf{Milestone:} independent-transfer-milestone\par
\noindent\textbf{Skill:} Revisit the skill on a fresh problem\par
\noindent\textbf{Experiment:} After a delay, attempt the fresh task and produce a reproducible evidence packet. Ask a peer to select or inspect the variation when possible; record assistance, failed hypotheses and remaining limits.\par
\noindent\textbf{Transfer:} Change a mechanism-relevant condition, rather than only resetting the old environment, and independently predict and test the resulting behavior.\par
\noindent\textbf{Control:} Include a contrasting case where the mechanism should not apply and a legitimate operation that should remain valid.\par
\noindent\textbf{Later retrieval:} Explain what transferred, what required help and which prerequisite needs another small experiment.\par
\noindent\textbf{Alignment references:} microsoft-ad, microsoft-adcs, goad, htb-pentest\par
\subsubsection{cloud-security: foundations}
\noindent\textbf{Milestone:} foundations-milestone\par
\noindent\textbf{Skill:} Prove the entry skills in a small lab\par
\noindent\textbf{Experiment:} Demonstrate each entry skill with a small owned fixture or permitted training environment. Record the actual setup, one observation per skill and any prerequisite you need to repair.\par
\noindent\textbf{Transfer:} Change a relevant identity, input or environment setting and predict the effect before repeating the demonstration.\par
\noindent\textbf{Control:} Include a deliberately unmet prerequisite and distinguish setup failure from a result about the subject.\par
\noindent\textbf{Later retrieval:} Reconstruct the environment and entry assumptions without following the setup notes.\par
\noindent\textbf{Alignment references:} owasp-cloud, cloudgoat, microsoft-soc\par
\subsubsection{cloud-security: shared-responsibility}
\noindent\textbf{Milestone:} shared-responsibility-milestone\par
\noindent\textbf{Skill:} Identify who owns each control\par
\noindent\textbf{Experiment:} Build an asset and responsibility inventory covering account ownership, regions, exposed endpoints, audit sources and recovery responsibilities.\par
\noindent\textbf{Transfer:} Move one component from a VM to a managed service and revise the shared-responsibility assumptions.\par
\noindent\textbf{Control:} Include a provider-managed control that a workload operator cannot fix by changing application code.\par
\noindent\textbf{Later retrieval:} Explain the difference between provider security and secure use of the provider.\par
\noindent\textbf{Alignment references:} owasp-cloud, cloudgoat, microsoft-soc\par
\subsubsection{cloud-security: iam}
\noindent\textbf{Milestone:} iam-milestone\par
\noindent\textbf{Skill:} Predict the authorization decision\par
\noindent\textbf{Experiment:} Build an access matrix for two synthetic roles and one resource, test it in a provider policy simulator or scoped training account, and record where the simulator lacks runtime context.\par
\noindent\textbf{Transfer:} Add a condition or a resource policy and predict how the effective decision changes.\par
\noindent\textbf{Control:} Demonstrate a denied unrelated principal and an intended allowed operation.\par
\noindent\textbf{Later retrieval:} Explain why an identity policy alone may not determine effective access.\par
\noindent\textbf{Alignment references:} owasp-cloud, cloudgoat, microsoft-soc\par
\subsubsection{cloud-security: exposure}
\noindent\textbf{Milestone:} exposure-milestone\par
\noindent\textbf{Skill:} Public does not always mean intended\par
\noindent\textbf{Experiment:} Trace a fictional application's exposure through load balancer, workload and data storage; validate selected rules using a local model or training account and propose a minimal repair.\par
\noindent\textbf{Transfer:} Change a resource from private to shared-with-one-role and verify that the intended boundary survives.\par
\noindent\textbf{Control:} Preserve legitimate client access and demonstrate that an unrelated identity cannot read the object.\par
\noindent\textbf{Later retrieval:} Reconstruct the layers required to explain one exposure finding.\par
\noindent\textbf{Alignment references:} owasp-cloud, cloudgoat, microsoft-soc\par
\subsubsection{cloud-security: secrets}
\noindent\textbf{Milestone:} secrets-milestone\par
\noindent\textbf{Skill:} Separate workload identity from static keys\par
\noindent\textbf{Experiment:} Replace a toy long-lived credential with a bounded local token or a training-provider workload identity, test expiration/revocation and document remaining trust assumptions.\par
\noindent\textbf{Transfer:} Change the audience, role or expiry and predict whether the request should still authenticate.\par
\noindent\textbf{Control:} Verify a valid authorized request succeeds and a stale token is rejected without logging its contents.\par
\noindent\textbf{Later retrieval:} Explain rotation, revocation and short-lived identity as different mechanisms.\par
\noindent\textbf{Alignment references:} owasp-cloud, cloudgoat, microsoft-soc\par
\subsubsection{cloud-security: logging}
\noindent\textbf{Milestone:} logging-milestone\par
\noindent\textbf{Skill:} Follow a control-plane change\par
\noindent\textbf{Experiment:} Investigate supplied cloud audit events around a permission change and later resource access; construct a timeline and a detection query with explicit missing sources.\par
\noindent\textbf{Transfer:} Change a resource name or actor and test whether the behavioral analytic still finds the sequence.\par
\noindent\textbf{Control:} Include a denied change and a known maintenance change that must not be labeled successful compromise.\par
\noindent\textbf{Later retrieval:} Explain what the audit source can and cannot establish about data access.\par
\noindent\textbf{Alignment references:} owasp-cloud, cloudgoat, microsoft-soc\par
\subsubsection{cloud-security: assessment}
\noindent\textbf{Milestone:} assessment-milestone\par
\noindent\textbf{Skill:} Deliver a bounded cloud review\par
\noindent\textbf{Experiment:} Review a fresh configuration and audit bundle, verify one identity or exposure defect in a scoped lab, apply a repair and record cleanup; use only owned training resources.\par
\noindent\textbf{Transfer:} Change one trust policy or network boundary and independently reassess the affected access path.\par
\noindent\textbf{Control:} Verify legitimate workload behavior and prove temporary lab resources were removed where a real account was used.\par
\noindent\textbf{Later retrieval:} Reconstruct the finding, provider-specific limits and remaining assurance gaps.\par
\noindent\textbf{Alignment references:} owasp-cloud, cloudgoat, microsoft-soc\par
\subsubsection{cloud-security: independent-transfer}
\noindent\textbf{Milestone:} independent-transfer-milestone\par
\noindent\textbf{Skill:} Revisit the skill on a fresh problem\par
\noindent\textbf{Experiment:} After a delay, attempt the fresh task and produce a reproducible evidence packet. Ask a peer to select or inspect the variation when possible; record assistance, failed hypotheses and remaining limits.\par
\noindent\textbf{Transfer:} Change a mechanism-relevant condition, rather than only resetting the old environment, and independently predict and test the resulting behavior.\par
\noindent\textbf{Control:} Include a contrasting case where the mechanism should not apply and a legitimate operation that should remain valid.\par
\noindent\textbf{Later retrieval:} Explain what transferred, what required help and which prerequisite needs another small experiment.\par
\noindent\textbf{Alignment references:} owasp-cloud, cloudgoat, microsoft-soc\par
\subsubsection{container-kubernetes-security: foundations}
\noindent\textbf{Milestone:} foundations-milestone\par
\noindent\textbf{Skill:} Prove the entry skills in a small lab\par
\noindent\textbf{Experiment:} Demonstrate each entry skill with a small owned fixture or permitted training environment. Record the actual setup, one observation per skill and any prerequisite you need to repair.\par
\noindent\textbf{Transfer:} Change a relevant identity, input or environment setting and predict the effect before repeating the demonstration.\par
\noindent\textbf{Control:} Include a deliberately unmet prerequisite and distinguish setup failure from a result about the subject.\par
\noindent\textbf{Later retrieval:} Reconstruct the environment and entry assumptions without following the setup notes.\par
\noindent\textbf{Alignment references:} kubernetes-checklist, owasp-cloud\par
\subsubsection{container-kubernetes-security: container-boundary}
\noindent\textbf{Milestone:} container-boundary-milestone\par
\noindent\textbf{Skill:} A container still shares a kernel\par
\noindent\textbf{Experiment:} Run a toy service with a non-root user, dropped capabilities, no-new-privileges and a read-only root plus a bounded writable directory; explain observed failures.\par
\noindent\textbf{Transfer:} Remove one required write path or change its owner and predict which operation fails.\par
\noindent\textbf{Control:} Preserve the service's intended request while showing it cannot write to an unrelated path.\par
\noindent\textbf{Later retrieval:} Explain which risks these controls reduce and which require stronger isolation.\par
\noindent\textbf{Alignment references:} kubernetes-checklist, owasp-cloud\par
\subsubsection{container-kubernetes-security: image-supply-chain}
\noindent\textbf{Milestone:} image-supply-chain-milestone\par
\noindent\textbf{Skill:} Know exactly which image ran\par
\noindent\textbf{Experiment:} Generate an image inventory/SBOM for a toy service, pin the source image digest and record provenance, vulnerability findings and applicability limits.\par
\noindent\textbf{Transfer:} Resolve an updated tag and identify the changed digest and dependencies before accepting it.\par
\noindent\textbf{Control:} Include a scanner finding for an unused or unreachable component and separate presence from demonstrated exploitability.\par
\noindent\textbf{Later retrieval:} Explain why a matching digest gives identity but not a guarantee of safety.\par
\noindent\textbf{Alignment references:} kubernetes-checklist, owasp-cloud\par
\subsubsection{container-kubernetes-security: network-separation}
\noindent\textbf{Milestone:} network-separation-milestone\par
\noindent\textbf{Skill:} Test an isolation claim\par
\noindent\textbf{Experiment:} In a disposable project, verify loopback-only access, separate lab networks, resource bounds and the absence of tracker-volume or Docker-socket mounts.\par
\noindent\textbf{Transfer:} Move the client to another lab network and predict which target it should fail to reach.\par
\noindent\textbf{Control:} Confirm a permitted same-network request works while cross-network and outbound probes fail within bounded timeouts.\par
\noindent\textbf{Later retrieval:} Explain what an internal bridge blocks and why the host kernel is still shared.\par
\noindent\textbf{Alignment references:} kubernetes-checklist, owasp-cloud\par
\subsubsection{container-kubernetes-security: cluster-identity}
\noindent\textbf{Milestone:} cluster-identity-milestone\par
\noindent\textbf{Skill:} Give the workload only its job\par
\noindent\textbf{Experiment:} In a disposable Kubernetes cluster, bind a limited service account to one namespace and verify allowed and denied operations with the API authorization checks.\par
\noindent\textbf{Transfer:} Change the namespace or requested verb and predict the decision before checking.\par
\noindent\textbf{Control:} Keep the intended read working and deny writes to unrelated secrets or namespaces.\par
\noindent\textbf{Later retrieval:} Reconstruct the subject, role and binding needed for one API permission.\par
\noindent\textbf{Alignment references:} kubernetes-checklist, owasp-cloud\par
\subsubsection{container-kubernetes-security: cluster-network}
\noindent\textbf{Milestone:} cluster-network-milestone\par
\noindent\textbf{Skill:} Observe enforcement, not just YAML\par
\noindent\textbf{Experiment:} Deploy a tiny client/server pair on an isolated test cluster, apply default-deny plus a specific allow rule, and test actual connectivity.\par
\noindent\textbf{Transfer:} Change one pod label or allowed port and predict which connection changes.\par
\noindent\textbf{Control:} Include a matching permitted client and an unmatched client that must be denied; report an unsupported CNI as an unmet prerequisite.\par
\noindent\textbf{Later retrieval:} Explain why a valid NetworkPolicy object can coexist with unrestricted traffic.\par
\noindent\textbf{Alignment references:} kubernetes-checklist, owasp-cloud\par
\subsubsection{container-kubernetes-security: review}
\noindent\textbf{Milestone:} review-milestone\par
\noindent\textbf{Skill:} Review and repair a workload deployment\par
\noindent\textbf{Experiment:} Review a fresh toy deployment, demonstrate one excessive permission or exposure, apply a minimal fix and retest the intended workload plus resource limits.\par
\noindent\textbf{Transfer:} Change a namespace, image version or writable path and revise the controls instead of copying the old manifest.\par
\noindent\textbf{Control:} Show an unrelated workload and legitimate client remain functional after the fix.\par
\noindent\textbf{Later retrieval:} Explain one runtime observation that would invalidate your review conclusion.\par
\noindent\textbf{Alignment references:} kubernetes-checklist, owasp-cloud\par
\subsubsection{container-kubernetes-security: independent-transfer}
\noindent\textbf{Milestone:} independent-transfer-milestone\par
\noindent\textbf{Skill:} Revisit the skill on a fresh problem\par
\noindent\textbf{Experiment:} After a delay, attempt the fresh task and produce a reproducible evidence packet. Ask a peer to select or inspect the variation when possible; record assistance, failed hypotheses and remaining limits.\par
\noindent\textbf{Transfer:} Change a mechanism-relevant condition, rather than only resetting the old environment, and independently predict and test the resulting behavior.\par
\noindent\textbf{Control:} Include a contrasting case where the mechanism should not apply and a legitimate operation that should remain valid.\par
\noindent\textbf{Later retrieval:} Explain what transferred, what required help and which prerequisite needs another small experiment.\par
\noindent\textbf{Alignment references:} kubernetes-checklist, owasp-cloud\par
\subsubsection{secure-development-devsecops: foundations}
\noindent\textbf{Milestone:} foundations-milestone\par
\noindent\textbf{Skill:} Prove the entry skills in a small lab\par
\noindent\textbf{Experiment:} Demonstrate each entry skill with a small owned fixture or permitted training environment. Record the actual setup, one observation per skill and any prerequisite you need to repair.\par
\noindent\textbf{Transfer:} Change a relevant identity, input or environment setting and predict the effect before repeating the demonstration.\par
\noindent\textbf{Control:} Include a deliberately unmet prerequisite and distinguish setup failure from a result about the subject.\par
\noindent\textbf{Later retrieval:} Reconstruct the environment and entry assumptions without following the setup notes.\par
\noindent\textbf{Alignment references:} owasp-asvs, owasp-cloud, kubernetes-checklist\par
\subsubsection{secure-development-devsecops: threat-model}
\noindent\textbf{Milestone:} threat-model-milestone\par
\noindent\textbf{Skill:} Draw the trust boundary before coding\par
\noindent\textbf{Experiment:} Build a threat model with three concrete abuse cases and map each to an ASVS-style testable requirement.\par
\noindent\textbf{Transfer:} Add a background worker or third-party integration and revise the boundaries and threats.\par
\noindent\textbf{Control:} Include a failure outside the service's responsibility and document the dependency instead of claiming to fix it.\par
\noindent\textbf{Later retrieval:} Reconstruct one threat from asset through entry point to consequence.\par
\noindent\textbf{Alignment references:} owasp-asvs, owasp-cloud, kubernetes-checklist\par
\subsubsection{secure-development-devsecops: input-and-output}
\noindent\textbf{Milestone:} input-and-output-milestone\par
\noindent\textbf{Skill:} Follow untrusted data to its sink\par
\noindent\textbf{Experiment:} Repair one input or output handling defect in your own code and write tests that reproduce the defect and preserve intended behavior.\par
\noindent\textbf{Transfer:} Change the input encoding or rendering context and update the defense rather than reusing an unrelated escaping rule.\par
\noindent\textbf{Control:} Include valid punctuation and an invalid structured input; reject the dangerous case without silently corrupting legitimate data.\par
\noindent\textbf{Later retrieval:} Explain why validation and output encoding address different boundaries.\par
\noindent\textbf{Alignment references:} owasp-asvs, owasp-cloud, kubernetes-checklist\par
\subsubsection{secure-development-devsecops: authz}
\noindent\textbf{Milestone:} authz-milestone\par
\noindent\textbf{Skill:} Put authorization beside the resource\par
\noindent\textbf{Experiment:} Implement server-side authorization tests for owner access, cross-user access and the privileged operation; demonstrate the bug before the repair.\par
\noindent\textbf{Transfer:} Add a new HTTP method or background operation and check that it enforces the same resource policy.\par
\noindent\textbf{Control:} Keep owner access working and deny an authenticated but unauthorized user.\par
\noindent\textbf{Later retrieval:} State the policy independently of route names or UI controls.\par
\noindent\textbf{Alignment references:} owasp-asvs, owasp-cloud, kubernetes-checklist\par
\subsubsection{secure-development-devsecops: supply-chain}
\noindent\textbf{Milestone:} supply-chain-milestone\par
\noindent\textbf{Skill:} Make a dependency decision you can defend\par
\noindent\textbf{Experiment:} Produce an SBOM and pinned dependency manifest, investigate one advisory's applicability, and test an update or documented mitigation.\par
\noindent\textbf{Transfer:} Introduce a version change in a branch and assess both security and compatibility effects before merging.\par
\noindent\textbf{Control:} Include an advisory that does not apply to the used component/version and explain the evidence.\par
\noindent\textbf{Later retrieval:} Explain why dependency presence, reachability and exploitability are separate claims.\par
\noindent\textbf{Alignment references:} owasp-asvs, owasp-cloud, kubernetes-checklist\par
\subsubsection{secure-development-devsecops: pipeline}
\noindent\textbf{Milestone:} pipeline-milestone\par
\noindent\textbf{Skill:} Make the pipeline stop for a real defect\par
\noindent\textbf{Experiment:} Add security regression tests, a synthetic-secret check and artifact identity recording; demonstrate that each relevant bad fixture fails the gate.\par
\noindent\textbf{Transfer:} Change a fixture's filename or encoding and test whether the intended check still applies.\par
\noindent\textbf{Control:} Confirm clean code builds successfully and no test prints secret-like fixture values into published logs.\par
\noindent\textbf{Later retrieval:} Explain the failure mode each gate detects and one it misses.\par
\noindent\textbf{Alignment references:} owasp-asvs, owasp-cloud, kubernetes-checklist\par
\subsubsection{secure-development-devsecops: release}
\noindent\textbf{Milestone:} release-milestone\par
\noindent\textbf{Skill:} Ship a small service with an assurance record\par
\noindent\textbf{Experiment:} Deliver the toy service with passing regression tests, reproducible build inputs, deployment permissions, rollback instructions and a documented residual risk.\par
\noindent\textbf{Transfer:} Change a dependency or deployment setting and rerun only the checks whose assumptions changed, explaining the choice.\par
\noindent\textbf{Control:} Reproduce a known security regression on the old version and show the release rejects it while valid use works.\par
\noindent\textbf{Later retrieval:} Reconstruct the assurance argument and the strongest remaining limitation.\par
\noindent\textbf{Alignment references:} owasp-asvs, owasp-cloud, kubernetes-checklist\par
\subsubsection{secure-development-devsecops: independent-transfer}
\noindent\textbf{Milestone:} independent-transfer-milestone\par
\noindent\textbf{Skill:} Revisit the skill on a fresh problem\par
\noindent\textbf{Experiment:} After a delay, attempt the fresh task and produce a reproducible evidence packet. Ask a peer to select or inspect the variation when possible; record assistance, failed hypotheses and remaining limits.\par
\noindent\textbf{Transfer:} Change a mechanism-relevant condition, rather than only resetting the old environment, and independently predict and test the resulting behavior.\par
\noindent\textbf{Control:} Include a contrasting case where the mechanism should not apply and a legitimate operation that should remain valid.\par
\noindent\textbf{Later retrieval:} Explain what transferred, what required help and which prerequisite needs another small experiment.\par
\noindent\textbf{Alignment references:} owasp-asvs, owasp-cloud, kubernetes-checklist\par
\subsubsection{digital-forensics-incident-response: foundations}
\noindent\textbf{Milestone:} foundations-milestone\par
\noindent\textbf{Skill:} Prove the entry skills in a small lab\par
\noindent\textbf{Experiment:} Demonstrate each entry skill with a small owned fixture or permitted training environment. Record the actual setup, one observation per skill and any prerequisite you need to repair.\par
\noindent\textbf{Transfer:} Change a relevant identity, input or environment setting and predict the effect before repeating the demonstration.\par
\noindent\textbf{Control:} Include a deliberately unmet prerequisite and distinguish setup failure from a result about the subject.\par
\noindent\textbf{Later retrieval:} Reconstruct the environment and entry assumptions without following the setup notes.\par
\noindent\textbf{Alignment references:} digitalcorpora, security-onion-import, htb-soc\par
\subsubsection{digital-forensics-incident-response: preservation}
\noindent\textbf{Milestone:} preservation-milestone\par
\noindent\textbf{Skill:} Keep the original evidence original\par
\noindent\textbf{Experiment:} Create an evidence manifest with acquisition time, source, hashes, handling history and analysis-copy paths; verify the original remains unchanged after analysis.\par
\noindent\textbf{Transfer:} Change the copy or filename and determine which identity assertions still hold.\par
\noindent\textbf{Control:} Include a corrupted copy that must fail integrity verification and a byte-identical renamed copy that must pass content verification.\par
\noindent\textbf{Later retrieval:} Explain integrity, provenance and chain of custody as different properties.\par
\noindent\textbf{Alignment references:} digitalcorpora, security-onion-import, htb-soc\par
\subsubsection{digital-forensics-incident-response: filesystem}
\noindent\textbf{Milestone:} filesystem-milestone\par
\noindent\textbf{Skill:} Recover a story from file metadata\par
\noindent\textbf{Experiment:} Build a bounded account of file creation/modification using an analysis copy, recording filesystem semantics and tool/version assumptions.\par
\noindent\textbf{Transfer:} Change the timezone display or introduce a copied file whose timestamps differ from its real acquisition time.\par
\noindent\textbf{Control:} Include metadata that cannot prove who created a file and explicitly withhold that conclusion.\par
\noindent\textbf{Later retrieval:} Explain why one timestamp does not prove an event happened at that exact time.\par
\noindent\textbf{Alignment references:} digitalcorpora, security-onion-import, htb-soc\par
\subsubsection{digital-forensics-incident-response: timeline}
\noindent\textbf{Milestone:} timeline-milestone\par
\noindent\textbf{Skill:} Reconcile clocks before drawing conclusions\par
\noindent\textbf{Experiment:} Correlate a small authentication, process and network evidence bundle into a timeline with confidence and source references.\par
\noindent\textbf{Transfer:} Add a clock-skew correction or delayed event and revise the ordering where justified.\par
\noindent\textbf{Control:} Keep an unrelated similarly named process out of the incident chain.\par
\noindent\textbf{Later retrieval:} Reconstruct which event ordering is certain and which depends on clock assumptions.\par
\noindent\textbf{Alignment references:} digitalcorpora, security-onion-import, htb-soc\par
\subsubsection{digital-forensics-incident-response: memory}
\noindent\textbf{Milestone:} memory-milestone\par
\noindent\textbf{Skill:} Separate a process artifact from a verdict\par
\noindent\textbf{Experiment:} On a suitable analysis host, inspect process and connection artifacts from a public training image, compare them with the provided ground truth where available, and record unsupported interpretations.\par
\noindent\textbf{Transfer:} Change the filter or investigate a second process and independently test the same hypothesis.\par
\noindent\textbf{Control:} Include a benign process with a suspicious-looking name and avoid treating its name alone as malware proof.\par
\noindent\textbf{Later retrieval:} Explain how an incompatible symbol profile can invalidate an apparently plausible result.\par
\noindent\textbf{Alignment references:} digitalcorpora, security-onion-import, htb-soc\par
\subsubsection{digital-forensics-incident-response: containment}
\noindent\textbf{Milestone:} containment-milestone\par
\noindent\textbf{Skill:} Preserve the question while containing the risk\par
\noindent\textbf{Experiment:} Rehearse a synthetic incident using a disposable service: preserve relevant logs, disable only the affected test identity and verify intended service behavior.\par
\noindent\textbf{Transfer:} Remove access to one evidence source and revise the plan without inventing missing facts.\par
\noindent\textbf{Control:} Preserve a legitimate user's access and document evidence lost because of the chosen intervention.\par
\noindent\textbf{Later retrieval:} Explain one containment tradeoff and how it changes later interpretation.\par
\noindent\textbf{Alignment references:} digitalcorpora, security-onion-import, htb-soc\par
\subsubsection{digital-forensics-incident-response: case}
\noindent\textbf{Milestone:} case-milestone\par
\noindent\textbf{Skill:} Deliver an evidence-led incident report\par
\noindent\textbf{Experiment:} Produce an evidence-indexed timeline, affected scope, alternative hypotheses, response decision and remaining uncertainty from an unfamiliar dataset.\par
\noindent\textbf{Transfer:} Introduce one contradictory artifact and revise the explanation while retaining an audit of the earlier hypothesis.\par
\noindent\textbf{Control:} Identify one tempting conclusion that the available evidence does not support.\par
\noindent\textbf{Later retrieval:} Reconstruct the case from the evidence index and independently recheck one claim.\par
\noindent\textbf{Alignment references:} digitalcorpora, security-onion-import, htb-soc\par
\subsubsection{digital-forensics-incident-response: independent-transfer}
\noindent\textbf{Milestone:} independent-transfer-milestone\par
\noindent\textbf{Skill:} Revisit the skill on a fresh problem\par
\noindent\textbf{Experiment:} After a delay, attempt the fresh task and produce a reproducible evidence packet. Ask a peer to select or inspect the variation when possible; record assistance, failed hypotheses and remaining limits.\par
\noindent\textbf{Transfer:} Change a mechanism-relevant condition, rather than only resetting the old environment, and independently predict and test the resulting behavior.\par
\noindent\textbf{Control:} Include a contrasting case where the mechanism should not apply and a legitimate operation that should remain valid.\par
\noindent\textbf{Later retrieval:} Explain what transferred, what required help and which prerequisite needs another small experiment.\par
\noindent\textbf{Alignment references:} digitalcorpora, security-onion-import, htb-soc\par
\subsubsection{threat-hunting-intelligence: foundations}
\noindent\textbf{Milestone:} foundations-milestone\par
\noindent\textbf{Skill:} Prove the entry skills in a small lab\par
\noindent\textbf{Experiment:} Demonstrate each entry skill with a small owned fixture or permitted training environment. Record the actual setup, one observation per skill and any prerequisite you need to repair.\par
\noindent\textbf{Transfer:} Change a relevant identity, input or environment setting and predict the effect before repeating the demonstration.\par
\noindent\textbf{Control:} Include a deliberately unmet prerequisite and distinguish setup failure from a result about the subject.\par
\noindent\textbf{Later retrieval:} Reconstruct the environment and entry assumptions without following the setup notes.\par
\noindent\textbf{Alignment references:} mitre-emulation, stix, microsoft-soc\par
\subsubsection{threat-hunting-intelligence: source-quality}
\noindent\textbf{Milestone:} source-quality-milestone\par
\noindent\textbf{Skill:} Put confidence beside the claim\par
\noindent\textbf{Experiment:} Produce a short intelligence note with source dates, collection limits, confidence and the concrete defender decision it informs.\par
\noindent\textbf{Transfer:} Add a conflicting source or newer report and explain which claims change.\par
\noindent\textbf{Control:} Include a widely repeated claim without primary evidence and keep it below corroborated facts.\par
\noindent\textbf{Later retrieval:} Explain how source reliability differs from confidence in an individual claim.\par
\noindent\textbf{Alignment references:} mitre-emulation, stix, microsoft-soc\par
\subsubsection{threat-hunting-intelligence: behavior}
\noindent\textbf{Milestone:} behavior-milestone\par
\noindent\textbf{Skill:} Turn a report into a huntable hypothesis\par
\noindent\textbf{Experiment:} Write an ATT\&CK-aligned hunt hypothesis, list needed fields and expected benign lookalikes, and define evidence that would falsify it.\par
\noindent\textbf{Transfer:} Change the assumed platform or available sensor and revise or reject the hypothesis.\par
\noindent\textbf{Control:} Show a technique mapping that is too broad to justify the proposed query.\par
\noindent\textbf{Later retrieval:} Reconstruct the hypothesis without using an indicator list.\par
\noindent\textbf{Alignment references:} mitre-emulation, stix, microsoft-soc\par
\subsubsection{threat-hunting-intelligence: queries}
\noindent\textbf{Milestone:} queries-milestone\par
\noindent\textbf{Skill:} Search for behavior through the noise\par
\noindent\textbf{Experiment:} Implement a reproducible hunt over synthetic telemetry and measure matches against labeled positive and benign events.\par
\noindent\textbf{Transfer:} Change filenames, addresses or usernames while preserving behavior and evaluate coverage.\par
\noindent\textbf{Control:} Include an ordinary administration sequence with similar individual events.\par
\noindent\textbf{Later retrieval:} Explain the false-positive mechanism of the original query.\par
\noindent\textbf{Alignment references:} mitre-emulation, stix, microsoft-soc\par
\subsubsection{threat-hunting-intelligence: context}
\noindent\textbf{Milestone:} context-milestone\par
\noindent\textbf{Skill:} Enrich without inventing attribution\par
\noindent\textbf{Experiment:} Build a small evidence graph from supplied indicators and public context, assigning dates, confidence and alternative explanations to each edge.\par
\noindent\textbf{Transfer:} Introduce a shared hosting address or recycled domain and remove unjustified relationships.\par
\noindent\textbf{Control:} Demonstrate that shared infrastructure alone does not establish a common operator.\par
\noindent\textbf{Later retrieval:} Explain one edge that is useful operationally but insufficient for attribution.\par
\noindent\textbf{Alignment references:} mitre-emulation, stix, microsoft-soc\par
\subsubsection{threat-hunting-intelligence: intelligence-format}
\noindent\textbf{Milestone:} intelligence-format-milestone\par
\noindent\textbf{Skill:} Exchange the meaning, not just the fields\par
\noindent\textbf{Experiment:} Create and validate a minimal STIX bundle for a fictional incident using synthetic indicators, confidence and source references.\par
\noindent\textbf{Transfer:} Update an observation or revoke an indicator and express the change without silently erasing its history.\par
\noindent\textbf{Control:} Include a malformed relationship or unsupported claim that validation/review must reject.\par
\noindent\textbf{Later retrieval:} Reconstruct the difference between an observable, an indicator and an incident claim.\par
\noindent\textbf{Alignment references:} mitre-emulation, stix, microsoft-soc\par
\subsubsection{threat-hunting-intelligence: hunt-report}
\noindent\textbf{Milestone:} hunt-report-milestone\par
\noindent\textbf{Skill:} Turn a hunt into an operational decision\par
\noindent\textbf{Experiment:} Conduct the hunt, document query versions, investigated matches, negative results and coverage limits, then propose a detection or data-collection improvement.\par
\noindent\textbf{Transfer:} Add a changed positive sequence or remove a source and revise the result transparently.\par
\noindent\textbf{Control:} Keep absence of evidence separate from proof that the behavior never occurred.\par
\noindent\textbf{Later retrieval:} Explain what action the hunt supports and what further collection is needed.\par
\noindent\textbf{Alignment references:} mitre-emulation, stix, microsoft-soc\par
\subsubsection{threat-hunting-intelligence: independent-transfer}
\noindent\textbf{Milestone:} independent-transfer-milestone\par
\noindent\textbf{Skill:} Revisit the skill on a fresh problem\par
\noindent\textbf{Experiment:} After a delay, attempt the fresh task and produce a reproducible evidence packet. Ask a peer to select or inspect the variation when possible; record assistance, failed hypotheses and remaining limits.\par
\noindent\textbf{Transfer:} Change a mechanism-relevant condition, rather than only resetting the old environment, and independently predict and test the resulting behavior.\par
\noindent\textbf{Control:} Include a contrasting case where the mechanism should not apply and a legitimate operation that should remain valid.\par
\noindent\textbf{Later retrieval:} Explain what transferred, what required help and which prerequisite needs another small experiment.\par
\noindent\textbf{Alignment references:} mitre-emulation, stix, microsoft-soc\par
\subsubsection{malware-analysis: foundations}
\noindent\textbf{Milestone:} foundations-milestone\par
\noindent\textbf{Skill:} Prove the entry skills in a small lab\par
\noindent\textbf{Experiment:} Demonstrate each entry skill with a small owned fixture or permitted training environment. Record the actual setup, one observation per skill and any prerequisite you need to repair.\par
\noindent\textbf{Transfer:} Change a relevant identity, input or environment setting and predict the effect before repeating the demonstration.\par
\noindent\textbf{Control:} Include a deliberately unmet prerequisite and distinguish setup failure from a result about the subject.\par
\noindent\textbf{Later retrieval:} Reconstruct the environment and entry assumptions without following the setup notes.\par
\noindent\textbf{Alignment references:} flare-vm, remnux, pwn-college\par
\subsubsection{malware-analysis: triage}
\noindent\textbf{Milestone:} triage-milestone\par
\noindent\textbf{Skill:} Describe the file before naming a family\par
\noindent\textbf{Experiment:} Produce a static triage sheet separating measured properties from inferred behavior and identify one next question for deeper analysis.\par
\noindent\textbf{Transfer:} Change compiler options or append harmless data and determine which identifiers and structural properties change.\par
\noindent\textbf{Control:} Include a packed-looking or high-entropy benign file and avoid classifying it as malicious solely from entropy.\par
\noindent\textbf{Later retrieval:} Explain what a hash, import and string can each establish.\par
\noindent\textbf{Alignment references:} flare-vm, remnux, pwn-college\par
\subsubsection{malware-analysis: static-flow}
\noindent\textbf{Milestone:} static-flow-milestone\par
\noindent\textbf{Skill:} Follow the branch that causes the effect\par
\noindent\textbf{Experiment:} Trace the condition, data flow and effect in source and disassembly; predict two inputs before running the benign program.\par
\noindent\textbf{Transfer:} Change the input or compile optimization and re-identify the controlling condition.\par
\noindent\textbf{Control:} Include an unreachable suspicious-looking string that does not prove execution.\par
\noindent\textbf{Later retrieval:} Reconstruct the condition-to-effect chain without searching for keywords.\par
\noindent\textbf{Alignment references:} flare-vm, remnux, pwn-college\par
\subsubsection{malware-analysis: behavior}
\noindent\textbf{Milestone:} behavior-milestone\par
\noindent\textbf{Skill:} Investigate a trace without executing a sample\par
\noindent\textbf{Experiment:} Build a behavior timeline from published or synthetic telemetry and distinguish observed effects from capabilities only suggested by static analysis.\par
\noindent\textbf{Transfer:} Add a failed operation or contradictory timestamp and revise the behavior claims.\par
\noindent\textbf{Control:} Include an attempted connection that never succeeded and avoid calling it confirmed data transfer.\par
\noindent\textbf{Later retrieval:} Explain one gap that would require a different sensor or controlled VM analysis.\par
\noindent\textbf{Alignment references:} flare-vm, remnux, pwn-college\par
\subsubsection{malware-analysis: configuration}
\noindent\textbf{Milestone:} configuration-milestone\par
\noindent\textbf{Skill:} Decode only what the evidence supports\par
\noindent\textbf{Experiment:} Write a bounded parser for a self-authored encoded configuration, validate recovered fields and preserve the original byte offsets and assumptions.\par
\noindent\textbf{Transfer:} Change key order, length or encoding and make the parser fail clearly or support the documented variation.\par
\noindent\textbf{Control:} Include a malformed/truncated fixture that must not yield a fabricated valid configuration.\par
\noindent\textbf{Later retrieval:} Explain why successful decoding is not equivalent to execution or attribution.\par
\noindent\textbf{Alignment references:} flare-vm, remnux, pwn-college\par
\subsubsection{malware-analysis: signatures}
\noindent\textbf{Milestone:} signatures-milestone\par
\noindent\textbf{Skill:} Write a rule that survives a renamed file\par
\noindent\textbf{Experiment:} Write a YARA or equivalent declarative rule for the training family and measure true/false matches on a documented sample set.\par
\noindent\textbf{Transfer:} Rename or recompile the toy sample and evaluate which features survive.\par
\noindent\textbf{Control:} Include unrelated binaries sharing a common string so an overbroad signature fails the assessment.\par
\noindent\textbf{Later retrieval:} Explain one expected false positive and how to investigate it.\par
\noindent\textbf{Alignment references:} flare-vm, remnux, pwn-college\par
\subsubsection{malware-analysis: report}
\noindent\textbf{Milestone:} report-milestone\par
\noindent\textbf{Skill:} Present a defensible analysis, not a label\par
\noindent\textbf{Experiment:} Produce a report with static evidence, trace evidence, configuration interpretation, a bounded detection rule and explicit uncertainty; do not execute unknown malware.\par
\noindent\textbf{Transfer:} Change one fixture property or trace event and independently revise the affected conclusion.\par
\noindent\textbf{Control:} Demonstrate a claim you deliberately withheld because the artifact or trace was insufficient.\par
\noindent\textbf{Later retrieval:} Reconstruct the strongest finding and the observation that could disprove it.\par
\noindent\textbf{Alignment references:} flare-vm, remnux, pwn-college\par
\subsubsection{malware-analysis: independent-transfer}
\noindent\textbf{Milestone:} independent-transfer-milestone\par
\noindent\textbf{Skill:} Revisit the skill on a fresh problem\par
\noindent\textbf{Experiment:} After a delay, attempt the fresh task and produce a reproducible evidence packet. Ask a peer to select or inspect the variation when possible; record assistance, failed hypotheses and remaining limits.\par
\noindent\textbf{Transfer:} Change a mechanism-relevant condition, rather than only resetting the old environment, and independently predict and test the resulting behavior.\par
\noindent\textbf{Control:} Include a contrasting case where the mechanism should not apply and a legitimate operation that should remain valid.\par
\noindent\textbf{Later retrieval:} Explain what transferred, what required help and which prerequisite needs another small experiment.\par
\noindent\textbf{Alignment references:} flare-vm, remnux, pwn-college\par
\subsubsection{mobile-application-security: foundations}
\noindent\textbf{Milestone:} foundations-milestone\par
\noindent\textbf{Skill:} Prove the entry skills in a small lab\par
\noindent\textbf{Experiment:} Demonstrate each entry skill with a small owned fixture or permitted training environment. Record the actual setup, one observation per skill and any prerequisite you need to repair.\par
\noindent\textbf{Transfer:} Change a relevant identity, input or environment setting and predict the effect before repeating the demonstration.\par
\noindent\textbf{Control:} Include a deliberately unmet prerequisite and distinguish setup failure from a result about the subject.\par
\noindent\textbf{Later retrieval:} Reconstruct the environment and entry assumptions without following the setup notes.\par
\noindent\textbf{Alignment references:} mastg, owasp-asvs\par
\subsubsection{mobile-application-security: platform-model}
\noindent\textbf{Milestone:} platform-model-milestone\par
\noindent\textbf{Skill:} Locate the mobile trust boundary\par
\noindent\textbf{Experiment:} Inventory an owned practice app's manifest/entitlements, components and data flows and map two questions to MASTG tests.\par
\noindent\textbf{Transfer:} Add an exported component or a new permission in a toy build and reassess exposure.\par
\noindent\textbf{Control:} Include a permission declaration that does not establish actual sensitive behavior by itself.\par
\noindent\textbf{Later retrieval:} Explain how platform permissions and backend authorization differ.\par
\noindent\textbf{Alignment references:} mastg, owasp-asvs\par
\subsubsection{mobile-application-security: local-storage}
\noindent\textbf{Milestone:} local-storage-milestone\par
\noindent\textbf{Skill:} Find where the app leaves data\par
\noindent\textbf{Experiment:} Using an owned device/emulator and a sanctioned practice app, locate a synthetic sensitive value, document its storage/access assumptions and test a safer toy implementation.\par
\noindent\textbf{Transfer:} Change device state or storage location and predict whether the same extraction method applies.\par
\noindent\textbf{Control:} Keep intended app persistence working and show that unrelated sandboxed apps do not gain access through the proposed fix.\par
\noindent\textbf{Later retrieval:} Explain encryption-at-rest limits when the app itself can access the key.\par
\noindent\textbf{Alignment references:} mastg, owasp-asvs\par
\subsubsection{mobile-application-security: transport}
\noindent\textbf{Milestone:} transport-milestone\par
\noindent\textbf{Skill:} Inspect a request without weakening the conclusion\par
\noindent\textbf{Experiment:} Capture permitted training-app traffic on your own device/emulator, identify transmitted synthetic data and evaluate server identity validation under documented test settings.\par
\noindent\textbf{Transfer:} Change the server certificate or hostname and predict acceptance or rejection.\par
\noindent\textbf{Control:} Verify a valid connection succeeds and an invalid identity is rejected by the corrected toy app.\par
\noindent\textbf{Later retrieval:} Explain why disabling pinning for analysis does not prove a production bypass.\par
\noindent\textbf{Alignment references:} mastg, owasp-asvs\par
\subsubsection{mobile-application-security: components}
\noindent\textbf{Milestone:} components-milestone\par
\noindent\textbf{Skill:} Test the component's actual caller\par
\noindent\textbf{Experiment:} Test a documented component boundary with benign input and record the caller identity, required permissions and resulting app behavior.\par
\noindent\textbf{Transfer:} Change a deep-link parameter or invoking app identity and predict whether access should be permitted.\par
\noindent\textbf{Control:} Keep the intended navigation working while rejecting an unauthorized or malformed invocation.\par
\noindent\textbf{Later retrieval:} Explain why an exported component is a review target rather than automatically a vulnerability.\par
\noindent\textbf{Alignment references:} mastg, owasp-asvs\par
\subsubsection{mobile-application-security: backend}
\noindent\textbf{Milestone:} backend-milestone\par
\noindent\textbf{Skill:} The phone is not the authorization boundary\par
\noindent\textbf{Experiment:} Exercise a training app's API with two roles and verify server-side ownership checks independently of the mobile interface.\par
\noindent\textbf{Transfer:} Change object ID or request method and repeat the authorization decision with fresh evidence.\par
\noindent\textbf{Control:} Preserve authorized access and deny an authenticated user who does not own the resource.\par
\noindent\textbf{Later retrieval:} Explain which finding belongs to the mobile client and which belongs to the backend.\par
\noindent\textbf{Alignment references:} mastg, owasp-asvs\par
\subsubsection{mobile-application-security: assessment}
\noindent\textbf{Milestone:} assessment-milestone\par
\noindent\textbf{Skill:} Deliver a platform-aware mobile review\par
\noindent\textbf{Experiment:} Review a fresh practice app against selected MASTG objectives, produce reproducible findings and retest one correction in an owned build.\par
\noindent\textbf{Transfer:} Change an OS/app version or component permission and independently re-evaluate the affected finding.\par
\noindent\textbf{Control:} Record a test whose result depends on debug/rooted-device conditions and limit the production claim accordingly.\par
\noindent\textbf{Later retrieval:} Reconstruct one finding including the exact platform assumptions required to reproduce it.\par
\noindent\textbf{Alignment references:} mastg, owasp-asvs\par
\subsubsection{mobile-application-security: independent-transfer}
\noindent\textbf{Milestone:} independent-transfer-milestone\par
\noindent\textbf{Skill:} Revisit the skill on a fresh problem\par
\noindent\textbf{Experiment:} After a delay, attempt the fresh task and produce a reproducible evidence packet. Ask a peer to select or inspect the variation when possible; record assistance, failed hypotheses and remaining limits.\par
\noindent\textbf{Transfer:} Change a mechanism-relevant condition, rather than only resetting the old environment, and independently predict and test the resulting behavior.\par
\noindent\textbf{Control:} Include a contrasting case where the mechanism should not apply and a legitimate operation that should remain valid.\par
\noindent\textbf{Later retrieval:} Explain what transferred, what required help and which prerequisite needs another small experiment.\par
\noindent\textbf{Alignment references:} mastg, owasp-asvs\par
\subsubsection{api-security: foundations}
\noindent\textbf{Milestone:} foundations-milestone\par
\noindent\textbf{Skill:} Prove the entry skills in a small lab\par
\noindent\textbf{Experiment:} Demonstrate each entry skill with a small owned fixture or permitted training environment. Record the actual setup, one observation per skill and any prerequisite you need to repair.\par
\noindent\textbf{Transfer:} Change a relevant identity, input or environment setting and predict the effect before repeating the demonstration.\par
\noindent\textbf{Control:} Include a deliberately unmet prerequisite and distinguish setup failure from a result about the subject.\par
\noindent\textbf{Later retrieval:} Reconstruct the environment and entry assumptions without following the setup notes.\par
\noindent\textbf{Alignment references:} owasp-api, portswigger-paths, owasp-asvs\par
\subsubsection{api-security: inventory}
\noindent\textbf{Milestone:} inventory-milestone\par
\noindent\textbf{Skill:} Find the contract the client assumes\par
\noindent\textbf{Experiment:} Inventory an owned training API, compare its documented and observed behavior, and identify two testable boundary hypotheses.\par
\noindent\textbf{Transfer:} Add a deprecated route or alternate version and update the inventory and risk assumptions.\par
\noindent\textbf{Control:} Distinguish an intentionally public endpoint from an unintentionally unprotected operation.\par
\noindent\textbf{Later retrieval:} Reconstruct one API trust boundary from method, identity and object.\par
\noindent\textbf{Alignment references:} owasp-api, portswigger-paths, owasp-asvs\par
\subsubsection{api-security: object-authorization}
\noindent\textbf{Milestone:} object-authorization-milestone\par
\noindent\textbf{Skill:} Change the object, keep the identity\par
\noindent\textbf{Experiment:} In a designated API lesson, test owner and cross-user access using recorded requests, explain the missing check and demonstrate a corrected local equivalent.\par
\noindent\textbf{Transfer:} Change the object type or HTTP verb and test whether the same boundary is enforced.\par
\noindent\textbf{Control:} Show legitimate owner access and rejection of an unrelated authenticated user.\par
\noindent\textbf{Later retrieval:} Explain why unpredictable IDs cannot replace authorization.\par
\noindent\textbf{Alignment references:} owasp-api, portswigger-paths, owasp-asvs\par
\subsubsection{api-security: property-authorization}
\noindent\textbf{Milestone:} property-authorization-milestone\par
\noindent\textbf{Skill:} Who may write this field?\par
\noindent\textbf{Experiment:} Test a local API's property allowlist with a harmless administrative-flag fixture and patch overbroad binding or response exposure.\par
\noindent\textbf{Transfer:} Nest the property or switch update methods and predict whether the policy should change.\par
\noindent\textbf{Control:} Keep ordinary editable fields working while rejecting privileged field writes and suppressing private response fields.\par
\noindent\textbf{Later retrieval:} Explain the relationship between mass assignment and field-level authorization.\par
\noindent\textbf{Alignment references:} owasp-api, portswigger-paths, owasp-asvs\par
\subsubsection{api-security: tokens}
\noindent\textbf{Milestone:} tokens-milestone\par
\noindent\textbf{Skill:} Validate the whole token contract\par
\noindent\textbf{Experiment:} Write or execute tests against an owned API for valid and invalid token contexts, separating signature validity from authorization.\par
\noindent\textbf{Transfer:} Change audience, scope or expiry and predict which validation layer rejects the request.\par
\noindent\textbf{Control:} Include a correctly signed but wrong-audience token and an authorized valid token.\par
\noindent\textbf{Later retrieval:} Explain why a valid signature does not make every action authorized.\par
\noindent\textbf{Alignment references:} owasp-api, portswigger-paths, owasp-asvs\par
\subsubsection{api-security: resource-limits}
\noindent\textbf{Milestone:} resource-limits-milestone\par
\noindent\textbf{Skill:} Make expensive work bounded\par
\noindent\textbf{Experiment:} Apply explicit request-size and work limits to an owned small service, then test bounded valid and over-limit inputs without load-testing third parties.\par
\noindent\textbf{Transfer:} Change nesting depth or batch composition and predict the work budget decision.\par
\noindent\textbf{Control:} Preserve a legitimate paginated request and reject an inexpensive-looking input that exceeds the actual work budget.\par
\noindent\textbf{Later retrieval:} Explain why request count alone may not bound resource consumption.\par
\noindent\textbf{Alignment references:} owasp-api, portswigger-paths, owasp-asvs\par
\subsubsection{api-security: review}
\noindent\textbf{Milestone:} review-milestone\par
\noindent\textbf{Skill:} Test a workflow across its endpoints\par
\noindent\textbf{Experiment:} Assess an unfamiliar owned API fixture for object, property, token and workflow boundaries; deliver evidence and a tested minimal repair.\par
\noindent\textbf{Transfer:} Change operation order or user role and independently test whether the business invariant still holds.\par
\noindent\textbf{Control:} Verify normal workflow completion and a rejected unauthorized transition after the fix.\par
\noindent\textbf{Later retrieval:} Reconstruct the finding from the business invariant rather than the endpoint name.\par
\noindent\textbf{Alignment references:} owasp-api, portswigger-paths, owasp-asvs\par
\subsubsection{api-security: independent-transfer}
\noindent\textbf{Milestone:} independent-transfer-milestone\par
\noindent\textbf{Skill:} Revisit the skill on a fresh problem\par
\noindent\textbf{Experiment:} After a delay, attempt the fresh task and produce a reproducible evidence packet. Ask a peer to select or inspect the variation when possible; record assistance, failed hypotheses and remaining limits.\par
\noindent\textbf{Transfer:} Change a mechanism-relevant condition, rather than only resetting the old environment, and independently predict and test the resulting behavior.\par
\noindent\textbf{Control:} Include a contrasting case where the mechanism should not apply and a legitimate operation that should remain valid.\par
\noindent\textbf{Later retrieval:} Explain what transferred, what required help and which prerequisite needs another small experiment.\par
\noindent\textbf{Alignment references:} owasp-api, portswigger-paths, owasp-asvs\par
\subsubsection{cryptography-protocol-security: foundations}
\noindent\textbf{Milestone:} foundations-milestone\par
\noindent\textbf{Skill:} Prove the entry skills in a small lab\par
\noindent\textbf{Experiment:} Demonstrate each entry skill with a small owned fixture or permitted training environment. Record the actual setup, one observation per skill and any prerequisite you need to repair.\par
\noindent\textbf{Transfer:} Change a relevant identity, input or environment setting and predict the effect before repeating the demonstration.\par
\noindent\textbf{Control:} Include a deliberately unmet prerequisite and distinguish setup failure from a result about the subject.\par
\noindent\textbf{Later retrieval:} Reconstruct the environment and entry assumptions without following the setup notes.\par
\noindent\textbf{Alignment references:} cryptohack, owasp-asvs\par
\subsubsection{cryptography-protocol-security: representations}
\noindent\textbf{Milestone:} representations-milestone\par
\noindent\textbf{Skill:} Bytes are not the text displaying them\par
\noindent\textbf{Experiment:} Implement round-trip tests for a bounded binary message and diagnose one encoding mismatch without losing bytes.\par
\noindent\textbf{Transfer:} Introduce a zero byte or non-ASCII character and predict the representation change.\par
\noindent\textbf{Control:} Reject malformed encoded input and preserve a valid binary message exactly.\par
\noindent\textbf{Later retrieval:} Explain why encoding provides no confidentiality.\par
\noindent\textbf{Alignment references:} cryptohack, owasp-asvs\par
\subsubsection{cryptography-protocol-security: integrity}
\noindent\textbf{Milestone:} integrity-milestone\par
\noindent\textbf{Skill:} A checksum does not authenticate its author\par
\noindent\textbf{Experiment:} Use a standard library to authenticate a message, test verification and record what the key ownership assumption means.\par
\noindent\textbf{Transfer:} Change one message byte or use a different key and predict verification.\par
\noindent\textbf{Control:} Accept the valid message and reject both tampering and a MAC made with the wrong key.\par
\noindent\textbf{Later retrieval:} Explain integrity versus origin authentication.\par
\noindent\textbf{Alignment references:} cryptohack, owasp-asvs\par
\subsubsection{cryptography-protocol-security: symmetric}
\noindent\textbf{Milestone:} symmetric-milestone\par
\noindent\textbf{Skill:} Make misuse visible in a toy experiment\par
\noindent\textbf{Experiment:} Encrypt and decrypt synthetic data with an established authenticated-encryption library and test tamper rejection.\par
\noindent\textbf{Transfer:} Change associated data or the nonce and predict which operation should fail.\par
\noindent\textbf{Control:} Keep the valid round trip working and reject a changed tag without releasing plaintext.\par
\noindent\textbf{Later retrieval:} Explain the nonce-use rule of the selected algorithm.\par
\noindent\textbf{Alignment references:} cryptohack, owasp-asvs\par
\subsubsection{cryptography-protocol-security: public-key}
\noindent\textbf{Milestone:} public-key-milestone\par
\noindent\textbf{Skill:} State exactly what a signature proves\par
\noindent\textbf{Experiment:} Sign and verify a synthetic document with an established library and connect the verified key to an explicit identity assumption.\par
\noindent\textbf{Transfer:} Change the message or trusted public key and predict the result.\par
\noindent\textbf{Control:} Reject a valid signature from an untrusted key while accepting the intended signer's document.\par
\noindent\textbf{Later retrieval:} Explain why a signature alone does not establish a real-world identity.\par
\noindent\textbf{Alignment references:} cryptohack, owasp-asvs\par
\subsubsection{cryptography-protocol-security: tls}
\noindent\textbf{Milestone:} tls-milestone\par
\noindent\textbf{Skill:} Test the certificate decision\par
\noindent\textbf{Experiment:} Run a local TLS test service with training certificates and verify hostname, trust-chain and expiry failures without disabling validation.\par
\noindent\textbf{Transfer:} Change the hostname or CA trust set and predict acceptance.\par
\noindent\textbf{Control:} Demonstrate a valid connection and rejection of an unrelated trusted-but-wrong-host certificate.\par
\noindent\textbf{Later retrieval:} Reconstruct the checks needed for server identity validation.\par
\noindent\textbf{Alignment references:} cryptohack, owasp-asvs\par
\subsubsection{cryptography-protocol-security: review}
\noindent\textbf{Milestone:} review-milestone\par
\noindent\textbf{Skill:} Find a protocol assumption that can be tested\par
\noindent\textbf{Experiment:} Review a toy protocol for replay, identity and integrity mistakes, repair one defect and provide executable regression tests.\par
\noindent\textbf{Transfer:} Reorder or replay a message and predict the corrected protocol's response.\par
\noindent\textbf{Control:} Preserve a legitimate fresh exchange while rejecting a stale authenticated message.\par
\noindent\textbf{Later retrieval:} Explain which property the repair establishes and which it does not.\par
\noindent\textbf{Alignment references:} cryptohack, owasp-asvs\par
\subsubsection{cryptography-protocol-security: independent-transfer}
\noindent\textbf{Milestone:} independent-transfer-milestone\par
\noindent\textbf{Skill:} Revisit the skill on a fresh problem\par
\noindent\textbf{Experiment:} After a delay, attempt the fresh task and produce a reproducible evidence packet. Ask a peer to select or inspect the variation when possible; record assistance, failed hypotheses and remaining limits.\par
\noindent\textbf{Transfer:} Change a mechanism-relevant condition, rather than only resetting the old environment, and independently predict and test the resulting behavior.\par
\noindent\textbf{Control:} Include a contrasting case where the mechanism should not apply and a legitimate operation that should remain valid.\par
\noindent\textbf{Later retrieval:} Explain what transferred, what required help and which prerequisite needs another small experiment.\par
\noindent\textbf{Alignment references:} cryptohack, owasp-asvs\par
\subsubsection{wireless-iot-security: foundations}
\noindent\textbf{Milestone:} foundations-milestone\par
\noindent\textbf{Skill:} Prove the entry skills in a small lab\par
\noindent\textbf{Experiment:} Demonstrate each entry skill with a small owned fixture or permitted training environment. Record the actual setup, one observation per skill and any prerequisite you need to repair.\par
\noindent\textbf{Transfer:} Change a relevant identity, input or environment setting and predict the effect before repeating the demonstration.\par
\noindent\textbf{Control:} Include a deliberately unmet prerequisite and distinguish setup failure from a result about the subject.\par
\noindent\textbf{Later retrieval:} Reconstruct the environment and entry assumptions without following the setup notes.\par
\noindent\textbf{Alignment references:} owasp-iot, mosquitto\par
\subsubsection{wireless-iot-security: inventory}
\noindent\textbf{Milestone:} inventory-milestone\par
\noindent\textbf{Skill:} Know which device and boundary you are testing\par
\noindent\textbf{Experiment:} Build a device threat model and an isolated test plan using an owned device or documented simulation.\par
\noindent\textbf{Transfer:} Add a companion cloud service and revise the data and trust boundaries.\par
\noindent\textbf{Control:} Exclude unrelated nearby devices even when technically visible.\par
\noindent\textbf{Later retrieval:} Explain which observations require physical hardware.\par
\noindent\textbf{Alignment references:} owasp-iot, mosquitto\par
\subsubsection{wireless-iot-security: mqtt}
\noindent\textbf{Milestone:} mqtt-milestone\par
\noindent\textbf{Skill:} Follow a message through the broker\par
\noindent\textbf{Experiment:} Publish and receive a unique harmless marker in the isolated MQTT profile and record exact topic/payload evidence.\par
\noindent\textbf{Transfer:} Change the subscription wildcard and predict which messages are delivered.\par
\noindent\textbf{Control:} Include an unmatched topic that must not reach the subscriber.\par
\noindent\textbf{Later retrieval:} Explain retained delivery versus a new live publication.\par
\noindent\textbf{Alignment references:} owasp-iot, mosquitto\par
\subsubsection{wireless-iot-security: access}
\noindent\textbf{Milestone:} access-milestone\par
\noindent\textbf{Skill:} Give each device only its topics\par
\noindent\textbf{Experiment:} Configure a separate owned broker with authentication and topic ACLs; verify each identity's allowed and denied operations.\par
\noindent\textbf{Transfer:} Change a wildcard or identity and predict effective topic access.\par
\noindent\textbf{Control:} Preserve a device's intended telemetry while denying another device's private topic.\par
\noindent\textbf{Later retrieval:} Reconstruct the difference between broker authentication and topic authorization.\par
\noindent\textbf{Alignment references:} owasp-iot, mosquitto\par
\subsubsection{wireless-iot-security: wireless}
\noindent\textbf{Milestone:} wireless-milestone\par
\noindent\textbf{Skill:} Interpret the capture you are actually able to make\par
\noindent\textbf{Experiment:} Analyze an owned or licensed training capture and explain association/security properties supported by visible frames.\par
\noindent\textbf{Transfer:} Change the selected channel or capture source and reassess missing observations.\par
\noindent\textbf{Control:} Avoid claiming to see encrypted application content without the required decryption context.\par
\noindent\textbf{Later retrieval:} Explain why a monitor-capable adapter is a prerequisite for some tests.\par
\noindent\textbf{Alignment references:} owasp-iot, mosquitto\par
\subsubsection{wireless-iot-security: firmware}
\noindent\textbf{Milestone:} firmware-milestone\par
\noindent\textbf{Skill:} Examine a firmware image as evidence\par
\noindent\textbf{Experiment:} Inspect static configuration and update metadata from a supplied image, identify one concrete trust assumption and propose a testable improvement.\par
\noindent\textbf{Transfer:} Change a signed metadata field or use a mismatched device model in a toy updater test.\par
\noindent\textbf{Control:} Accept valid matching metadata and reject tampering or a wrong-model update.\par
\noindent\textbf{Later retrieval:} Explain what static image analysis cannot prove about a running device.\par
\noindent\textbf{Alignment references:} owasp-iot, mosquitto\par
\subsubsection{wireless-iot-security: assessment}
\noindent\textbf{Milestone:} assessment-milestone\par
\noindent\textbf{Skill:} Review an end-to-end device workflow\par
\noindent\textbf{Experiment:} Assess a small owned MQTT workflow, repair one access defect and report remaining hardware/firmware assumptions.\par
\noindent\textbf{Transfer:} Introduce another device identity or a changed topic layout and retest separation.\par
\noindent\textbf{Control:} Keep legitimate telemetry working and deny the unrelated device.\par
\noindent\textbf{Later retrieval:} Reconstruct the finding through the complete device-to-service path.\par
\noindent\textbf{Alignment references:} owasp-iot, mosquitto\par
\subsubsection{wireless-iot-security: independent-transfer}
\noindent\textbf{Milestone:} independent-transfer-milestone\par
\noindent\textbf{Skill:} Revisit the skill on a fresh problem\par
\noindent\textbf{Experiment:} After a delay, attempt the fresh task and produce a reproducible evidence packet. Ask a peer to select or inspect the variation when possible; record assistance, failed hypotheses and remaining limits.\par
\noindent\textbf{Transfer:} Change a mechanism-relevant condition, rather than only resetting the old environment, and independently predict and test the resulting behavior.\par
\noindent\textbf{Control:} Include a contrasting case where the mechanism should not apply and a legitimate operation that should remain valid.\par
\noindent\textbf{Later retrieval:} Explain what transferred, what required help and which prerequisite needs another small experiment.\par
\noindent\textbf{Alignment references:} owasp-iot, mosquitto\par
\subsubsection{ot-ics-security: foundations}
\noindent\textbf{Milestone:} foundations-milestone\par
\noindent\textbf{Skill:} Prove the entry skills in a small lab\par
\noindent\textbf{Experiment:} Demonstrate each entry skill with a small owned fixture or permitted training environment. Record the actual setup, one observation per skill and any prerequisite you need to repair.\par
\noindent\textbf{Transfer:} Change a relevant identity, input or environment setting and predict the effect before repeating the demonstration.\par
\noindent\textbf{Control:} Include a deliberately unmet prerequisite and distinguish setup failure from a result about the subject.\par
\noindent\textbf{Later retrieval:} Reconstruct the environment and entry assumptions without following the setup notes.\par
\noindent\textbf{Alignment references:} cisa-ics, pymodbus, nist-csf\par
\subsubsection{ot-ics-security: process}
\noindent\textbf{Milestone:} process-milestone\par
\noindent\textbf{Skill:} Start with the process consequence\par
\noindent\textbf{Experiment:} Build an asset/zone diagram and identify how loss of view, control or availability affects the simulated process.\par
\noindent\textbf{Transfer:} Change a process dependency and revise the priority of affected controls.\par
\noindent\textbf{Control:} Distinguish a confidentiality issue from a change that could affect safe operation.\par
\noindent\textbf{Later retrieval:} Explain why enterprise IT priorities cannot be copied blindly into OT.\par
\noindent\textbf{Alignment references:} cisa-ics, pymodbus, nist-csf\par
\subsubsection{ot-ics-security: protocol}
\noindent\textbf{Milestone:} protocol-milestone\par
\noindent\textbf{Skill:} Read the register with its context\par
\noindent\textbf{Experiment:} Use an isolated software simulator or capture to trace a benign register read and document address and data-type assumptions.\par
\noindent\textbf{Transfer:} Change register address or byte order in a fixture and predict the interpretation.\par
\noindent\textbf{Control:} Reject an out-of-range or mismatched response instead of inventing a process value.\par
\noindent\textbf{Later retrieval:} Explain why a register number alone does not identify a physical measurement.\par
\noindent\textbf{Alignment references:} cisa-ics, pymodbus, nist-csf\par
\subsubsection{ot-ics-security: segmentation}
\noindent\textbf{Milestone:} segmentation-milestone\par
\noindent\textbf{Skill:} Prove the permitted path\par
\noindent\textbf{Experiment:} Model and test a small allowlist of communication paths using local fixtures or a disposable network.\par
\noindent\textbf{Transfer:} Change a client role or destination port and predict whether access should remain.\par
\noindent\textbf{Control:} Preserve required monitoring while denying an unrelated management path.\par
\noindent\textbf{Later retrieval:} Reconstruct one zone boundary and its operational justification.\par
\noindent\textbf{Alignment references:} cisa-ics, pymodbus, nist-csf\par
\subsubsection{ot-ics-security: monitoring}
\noindent\textbf{Milestone:} monitoring-milestone\par
\noindent\textbf{Skill:} Distinguish a maintenance action from an anomaly\par
\noindent\textbf{Experiment:} Create a detection hypothesis for an unexpected write or new peer and evaluate it against labeled simulator traffic.\par
\noindent\textbf{Transfer:} Change a nonessential address or schedule and reassess the rule.\par
\noindent\textbf{Control:} Include approved maintenance and record the context needed to avoid a false incident.\par
\noindent\textbf{Later retrieval:} Explain what the network sensor cannot establish about physical state.\par
\noindent\textbf{Alignment references:} cisa-ics, pymodbus, nist-csf\par
\subsubsection{ot-ics-security: change-control}
\noindent\textbf{Milestone:} change-control-milestone\par
\noindent\textbf{Skill:} Rehearse a reversible change\par
\noindent\textbf{Experiment:} Apply one non-disruptive simulator hardening change, preserve the baseline and verify expected readings afterward.\par
\noindent\textbf{Transfer:} Remove an assumed maintenance window and decide whether to defer or redesign the change.\par
\noindent\textbf{Control:} Demonstrate that a failed validation leads to rollback rather than an unsupported success claim.\par
\noindent\textbf{Later retrieval:} Explain the evidence required before a change could be considered in a real plant.\par
\noindent\textbf{Alignment references:} cisa-ics, pymodbus, nist-csf\par
\subsubsection{ot-ics-security: exercise}
\noindent\textbf{Milestone:} exercise-milestone\par
\noindent\textbf{Skill:} Run an incident tabletop with technical evidence\par
\noindent\textbf{Experiment:} Investigate a supplied OT trace and conduct a tabletop containment decision with a technical timeline and explicit safety uncertainty.\par
\noindent\textbf{Transfer:} Add contradictory process evidence and revise the decision.\par
\noindent\textbf{Control:} Reject an IT-style containment action that would violate the stated process constraint.\par
\noindent\textbf{Later retrieval:} Reconstruct the decision and the unresolved question requiring a process engineer.\par
\noindent\textbf{Alignment references:} cisa-ics, pymodbus, nist-csf\par
\subsubsection{ot-ics-security: independent-transfer}
\noindent\textbf{Milestone:} independent-transfer-milestone\par
\noindent\textbf{Skill:} Revisit the skill on a fresh problem\par
\noindent\textbf{Experiment:} After a delay, attempt the fresh task and produce a reproducible evidence packet. Ask a peer to select or inspect the variation when possible; record assistance, failed hypotheses and remaining limits.\par
\noindent\textbf{Transfer:} Change a mechanism-relevant condition, rather than only resetting the old environment, and independently predict and test the resulting behavior.\par
\noindent\textbf{Control:} Include a contrasting case where the mechanism should not apply and a legitimate operation that should remain valid.\par
\noindent\textbf{Later retrieval:} Explain what transferred, what required help and which prerequisite needs another small experiment.\par
\noindent\textbf{Alignment references:} cisa-ics, pymodbus, nist-csf\par
\subsubsection{security-architecture-grc: foundations}
\noindent\textbf{Milestone:} foundations-milestone\par
\noindent\textbf{Skill:} Prove the entry skills in a small lab\par
\noindent\textbf{Experiment:} Demonstrate each entry skill with a small owned fixture or permitted training environment. Record the actual setup, one observation per skill and any prerequisite you need to repair.\par
\noindent\textbf{Transfer:} Change a relevant identity, input or environment setting and predict the effect before repeating the demonstration.\par
\noindent\textbf{Control:} Include a deliberately unmet prerequisite and distinguish setup failure from a result about the subject.\par
\noindent\textbf{Later retrieval:} Reconstruct the environment and entry assumptions without following the setup notes.\par
\noindent\textbf{Alignment references:} nist-csf, owasp-asvs, cisa-ssvc\par
\subsubsection{security-architecture-grc: assets}
\noindent\textbf{Milestone:} assets-milestone\par
\noindent\textbf{Skill:} Define the system someone is responsible for\par
\noindent\textbf{Experiment:} Write a system boundary and data-flow diagram that includes outsourced services and critical business assumptions.\par
\noindent\textbf{Transfer:} Add a new data recipient and revise ownership and control responsibilities.\par
\noindent\textbf{Control:} Identify an asset outside the boundary that still creates a dependency risk.\par
\noindent\textbf{Later retrieval:} Explain why an incomplete boundary undermines an assessment.\par
\noindent\textbf{Alignment references:} nist-csf, owasp-asvs, cisa-ssvc\par
\subsubsection{security-architecture-grc: threats}
\noindent\textbf{Milestone:} threats-milestone\par
\noindent\textbf{Skill:} Connect the threat to a business consequence\par
\noindent\textbf{Experiment:} Create a small risk register with evidence, uncertainty, existing controls and a named decision owner.\par
\noindent\textbf{Transfer:} Change exposure or business criticality and justify the resulting priority.\par
\noindent\textbf{Control:} Avoid turning an unsupported numeric guess into a precise probability claim.\par
\noindent\textbf{Later retrieval:} Reconstruct one risk from cause to consequence.\par
\noindent\textbf{Alignment references:} nist-csf, owasp-asvs, cisa-ssvc\par
\subsubsection{security-architecture-grc: controls}
\noindent\textbf{Milestone:} controls-milestone\par
\noindent\textbf{Skill:} Make the control testable\par
\noindent\textbf{Experiment:} Map selected NIST CSF outcomes to concrete controls and evidence collection for a fictional service.\par
\noindent\textbf{Transfer:} Change the hosting model and identify controls whose evidence must change.\par
\noindent\textbf{Control:} Include a policy document that does not prove the operational control is running.\par
\noindent\textbf{Later retrieval:} Explain control design versus control operating effectiveness.\par
\noindent\textbf{Alignment references:} nist-csf, owasp-asvs, cisa-ssvc\par
\subsubsection{security-architecture-grc: architecture}
\noindent\textbf{Milestone:} architecture-milestone\par
\noindent\textbf{Skill:} Compare designs against explicit assumptions\par
\noindent\textbf{Experiment:} Evaluate the designs against the same threat model, operational needs and recovery requirements, then justify a bounded choice.\par
\noindent\textbf{Transfer:} Remove one assumed trusted component and revise the recommendation.\par
\noindent\textbf{Control:} Document one risk accepted by the chosen design instead of claiming universal superiority.\par
\noindent\textbf{Later retrieval:} Reconstruct the tradeoff that drove the decision.\par
\noindent\textbf{Alignment references:} nist-csf, owasp-asvs, cisa-ssvc\par
\subsubsection{security-architecture-grc: assurance}
\noindent\textbf{Milestone:} assurance-milestone\par
\noindent\textbf{Skill:} Trace a requirement to its evidence\par
\noindent\textbf{Experiment:} Review a small evidence pack, identify stale or missing evidence and produce a remediation plan with owners and retest criteria.\par
\noindent\textbf{Transfer:} Change a requirement version or evidence date and reassess adequacy.\par
\noindent\textbf{Control:} Reject a screenshot or scanner summary that cannot substantiate the actual claim.\par
\noindent\textbf{Later retrieval:} Explain the difference between a failed control and insufficient evidence.\par
\noindent\textbf{Alignment references:} nist-csf, owasp-asvs, cisa-ssvc\par
\subsubsection{security-architecture-grc: review}
\noindent\textbf{Milestone:} review-milestone\par
\noindent\textbf{Skill:} Deliver a decision-ready security review\par
\noindent\textbf{Experiment:} Deliver a concise architecture/risk review with prioritized actions, residual risks, evidence references and an accountable acceptance decision.\par
\noindent\textbf{Transfer:} Introduce a new dependency or material finding and update the decision transparently.\par
\noindent\textbf{Control:} Keep a low-impact observation from obscuring an unresolved launch-critical control.\par
\noindent\textbf{Later retrieval:} Reconstruct the strongest assurance claim and its limits.\par
\noindent\textbf{Alignment references:} nist-csf, owasp-asvs, cisa-ssvc\par
\subsubsection{security-architecture-grc: independent-transfer}
\noindent\textbf{Milestone:} independent-transfer-milestone\par
\noindent\textbf{Skill:} Revisit the skill on a fresh problem\par
\noindent\textbf{Experiment:} After a delay, attempt the fresh task and produce a reproducible evidence packet. Ask a peer to select or inspect the variation when possible; record assistance, failed hypotheses and remaining limits.\par
\noindent\textbf{Transfer:} Change a mechanism-relevant condition, rather than only resetting the old environment, and independently predict and test the resulting behavior.\par
\noindent\textbf{Control:} Include a contrasting case where the mechanism should not apply and a legitimate operation that should remain valid.\par
\noindent\textbf{Later retrieval:} Explain what transferred, what required help and which prerequisite needs another small experiment.\par
\noindent\textbf{Alignment references:} nist-csf, owasp-asvs, cisa-ssvc\par
\subsubsection{vulnerability-management: foundations}
\noindent\textbf{Milestone:} foundations-milestone\par
\noindent\textbf{Skill:} Prove the entry skills in a small lab\par
\noindent\textbf{Experiment:} Demonstrate each entry skill with a small owned fixture or permitted training environment. Record the actual setup, one observation per skill and any prerequisite you need to repair.\par
\noindent\textbf{Transfer:} Change a relevant identity, input or environment setting and predict the effect before repeating the demonstration.\par
\noindent\textbf{Control:} Include a deliberately unmet prerequisite and distinguish setup failure from a result about the subject.\par
\noindent\textbf{Later retrieval:} Reconstruct the environment and entry assumptions without following the setup notes.\par
\noindent\textbf{Alignment references:} cisa-kev, epss, cisa-ssvc\par
\subsubsection{vulnerability-management: inventory}
\noindent\textbf{Milestone:} inventory-milestone\par
\noindent\textbf{Skill:} Find the asset behind the finding\par
\noindent\textbf{Experiment:} Reconcile a synthetic scan inventory with deployed components and identify unknown ownership or stale records.\par
\noindent\textbf{Transfer:} Rename a host or move a service and preserve its identity without merging unrelated assets.\par
\noindent\textbf{Control:} Include two similar version strings that refer to different products.\par
\noindent\textbf{Later retrieval:} Explain why asset identity must precede prioritization.\par
\noindent\textbf{Alignment references:} cisa-kev, epss, cisa-ssvc\par
\subsubsection{vulnerability-management: validation}
\noindent\textbf{Milestone:} validation-milestone\par
\noindent\textbf{Skill:} Is the advisory actually applicable?\par
\noindent\textbf{Experiment:} Validate a supplied finding against exact component/configuration evidence and classify confirmed, not applicable or insufficient evidence.\par
\noindent\textbf{Transfer:} Change a required feature or version and revise applicability.\par
\noindent\textbf{Control:} Include a matching product name with a nonaffected version.\par
\noindent\textbf{Later retrieval:} Reconstruct the conditions required before calling the finding applicable.\par
\noindent\textbf{Alignment references:} cisa-kev, epss, cisa-ssvc\par
\subsubsection{vulnerability-management: prioritization}
\noindent\textbf{Milestone:} prioritization-milestone\par
\noindent\textbf{Skill:} Prioritize action using context\par
\noindent\textbf{Experiment:} Prioritize a synthetic queue using exposure, business impact, exploitation evidence and an explicit SSVC-style decision path.\par
\noindent\textbf{Transfer:} Change exposure or exploitation evidence and justify the revised action.\par
\noindent\textbf{Control:} Include a high-score isolated lab finding and a lower-score exposed critical asset.\par
\noindent\textbf{Later retrieval:} Explain which new evidence would change the highest priority.\par
\noindent\textbf{Alignment references:} cisa-kev, epss, cisa-ssvc\par
\subsubsection{vulnerability-management: remediation}
\noindent\textbf{Milestone:} remediation-milestone\par
\noindent\textbf{Skill:} Verify the repair rather than the ticket\par
\noindent\textbf{Experiment:} Apply a patch or configuration mitigation in an owned test service and verify both the vulnerable condition and normal functionality.\par
\noindent\textbf{Transfer:} Change the deployment variant and reassess whether the same mitigation applies.\par
\noindent\textbf{Control:} Keep a ticket marked resolved from passing when the affected instance still has the old configuration.\par
\noindent\textbf{Later retrieval:} Explain patch presence versus effective remediation.\par
\noindent\textbf{Alignment references:} cisa-kev, epss, cisa-ssvc\par
\subsubsection{vulnerability-management: exceptions}
\noindent\textbf{Milestone:} exceptions-milestone\par
\noindent\textbf{Skill:} Put an expiry date on accepted risk\par
\noindent\textbf{Experiment:} Write a bounded exception for a fictional unpatchable service, with measurable compensating controls and a retest trigger.\par
\noindent\textbf{Transfer:} Remove a compensating control or change external exposure and reevaluate the exception.\par
\noindent\textbf{Control:} Reject an indefinite exception without an accountable owner.\par
\noindent\textbf{Later retrieval:} Reconstruct when the exception must return to the remediation queue.\par
\noindent\textbf{Alignment references:} cisa-kev, epss, cisa-ssvc\par
\subsubsection{vulnerability-management: program}
\noindent\textbf{Milestone:} program-milestone\par
\noindent\textbf{Skill:} Measure reduction in exposure\par
\noindent\textbf{Experiment:} Produce a vulnerability-management report from a synthetic inventory, separating verified fixes, accepted risk, missing data and recurring causes.\par
\noindent\textbf{Transfer:} Add previously unknown assets and revise the metrics without hiding the changed denominator.\par
\noindent\textbf{Control:} Distinguish a fall in scan coverage from a real reduction in exposure.\par
\noindent\textbf{Later retrieval:} Explain one metric that could look better while security became worse.\par
\noindent\textbf{Alignment references:} cisa-kev, epss, cisa-ssvc\par
\subsubsection{vulnerability-management: independent-transfer}
\noindent\textbf{Milestone:} independent-transfer-milestone\par
\noindent\textbf{Skill:} Revisit the skill on a fresh problem\par
\noindent\textbf{Experiment:} After a delay, attempt the fresh task and produce a reproducible evidence packet. Ask a peer to select or inspect the variation when possible; record assistance, failed hypotheses and remaining limits.\par
\noindent\textbf{Transfer:} Change a mechanism-relevant condition, rather than only resetting the old environment, and independently predict and test the resulting behavior.\par
\noindent\textbf{Control:} Include a contrasting case where the mechanism should not apply and a legitimate operation that should remain valid.\par
\noindent\textbf{Later retrieval:} Explain what transferred, what required help and which prerequisite needs another small experiment.\par
\noindent\textbf{Alignment references:} cisa-kev, epss, cisa-ssvc\par
\subsubsection{vulnerability-research-bug-bounty: foundations}
\noindent\textbf{Milestone:} foundations-milestone\par
\noindent\textbf{Skill:} Prove the entry skills in a small lab\par
\noindent\textbf{Experiment:} Demonstrate each entry skill with a small owned fixture or permitted training environment. Record the actual setup, one observation per skill and any prerequisite you need to repair.\par
\noindent\textbf{Transfer:} Change a relevant identity, input or environment setting and predict the effect before repeating the demonstration.\par
\noindent\textbf{Control:} Include a deliberately unmet prerequisite and distinguish setup failure from a result about the subject.\par
\noindent\textbf{Later retrieval:} Reconstruct the environment and entry assumptions without following the setup notes.\par
\noindent\textbf{Alignment references:} portswigger-mystery, kettle-research, owasp-asvs\par
\subsubsection{vulnerability-research-bug-bounty: scope}
\noindent\textbf{Milestone:} scope-milestone\par
\noindent\textbf{Skill:} Treat program rules as part of the target\par
\noindent\textbf{Experiment:} Build a research plan using only an owned training target and a clearly testable vulnerability hypothesis.\par
\noindent\textbf{Transfer:} Add an excluded third-party integration and revise the plan.\par
\noindent\textbf{Control:} Identify a technically possible test that the scope disallows.\par
\noindent\textbf{Later retrieval:} Explain why a public bug-bounty page does not authorize every asset or technique.\par
\noindent\textbf{Alignment references:} portswigger-mystery, kettle-research, owasp-asvs\par
\subsubsection{vulnerability-research-bug-bounty: model}
\noindent\textbf{Milestone:} model-milestone\par
\noindent\textbf{Skill:} Find the invariant the software should preserve\par
\noindent\textbf{Experiment:} Form a root-cause hypothesis from the code and design a minimal experiment that could disprove it.\par
\noindent\textbf{Transfer:} Change input encoding or operation order and predict the effect on the invariant.\par
\noindent\textbf{Control:} Include a normal edge case that looks unusual but preserves the intended rule.\par
\noindent\textbf{Later retrieval:} Reconstruct the vulnerability hypothesis without relying on a scanner label.\par
\noindent\textbf{Alignment references:} portswigger-mystery, kettle-research, owasp-asvs\par
\subsubsection{vulnerability-research-bug-bounty: fuzzing}
\noindent\textbf{Milestone:} fuzzing-milestone\par
\noindent\textbf{Skill:} Make a small test generator teach you something\par
\noindent\textbf{Experiment:} Write a bounded local input generator or property-based test, reproduce one failure and minimize the triggering fixture.\par
\noindent\textbf{Transfer:} Change the parser or seed corpus and verify the failure still tests the same property.\par
\noindent\textbf{Control:} Include malformed input that is correctly rejected and must not count as a vulnerability.\par
\noindent\textbf{Later retrieval:} Explain crash, bug, security impact and exploitability as separate claims.\par
\noindent\textbf{Alignment references:} portswigger-mystery, kettle-research, owasp-asvs\par
\subsubsection{vulnerability-research-bug-bounty: impact}
\noindent\textbf{Milestone:} impact-milestone\par
\noindent\textbf{Skill:} Prove only the minimum necessary effect\par
\noindent\textbf{Experiment:} Demonstrate a harmless synthetic impact in an owned reproduction and document exactly which privileges and data were involved.\par
\noindent\textbf{Transfer:} Change the user role or object class and test whether the impact generalizes.\par
\noindent\textbf{Control:} Include a failed unauthorized action and avoid claiming access merely from a different response size.\par
\noindent\textbf{Later retrieval:} Explain one stronger impact claim the evidence does not support.\par
\noindent\textbf{Alignment references:} portswigger-mystery, kettle-research, owasp-asvs\par
\subsubsection{vulnerability-research-bug-bounty: disclosure}
\noindent\textbf{Milestone:} disclosure-milestone\par
\noindent\textbf{Skill:} Write a report that survives first triage\par
\noindent\textbf{Experiment:} Produce a private training report with minimal reproducer, sanitized evidence and a proposed tested fix; do not submit to third parties as part of this exercise.\par
\noindent\textbf{Transfer:} Have a peer or fresh environment follow the report with a changed fixture and repair any missing assumption.\par
\noindent\textbf{Control:} Keep credentials and live user data out of the evidence packet.\par
\noindent\textbf{Later retrieval:} Reconstruct the report's root cause in two precise sentences.\par
\noindent\textbf{Alignment references:} portswigger-mystery, kettle-research, owasp-asvs\par
\subsubsection{vulnerability-research-bug-bounty: research}
\noindent\textbf{Milestone:} research-milestone\par
\noindent\textbf{Skill:} Generalize the mechanism carefully\par
\noindent\textbf{Experiment:} Test the hypothesis on the new component, report positive or negative evidence, and produce a regression test and bounded generalization.\par
\noindent\textbf{Transfer:} Change one relevant implementation detail and explain why the mechanism applies or fails.\par
\noindent\textbf{Control:} Include a superficially similar implementation where the required condition is absent.\par
\noindent\textbf{Later retrieval:} Explain which properties transfer and which were accidents of the first target.\par
\noindent\textbf{Alignment references:} portswigger-mystery, kettle-research, owasp-asvs\par
\subsubsection{vulnerability-research-bug-bounty: independent-transfer}
\noindent\textbf{Milestone:} independent-transfer-milestone\par
\noindent\textbf{Skill:} Revisit the skill on a fresh problem\par
\noindent\textbf{Experiment:} After a delay, attempt the fresh task and produce a reproducible evidence packet. Ask a peer to select or inspect the variation when possible; record assistance, failed hypotheses and remaining limits.\par
\noindent\textbf{Transfer:} Change a mechanism-relevant condition, rather than only resetting the old environment, and independently predict and test the resulting behavior.\par
\noindent\textbf{Control:} Include a contrasting case where the mechanism should not apply and a legitimate operation that should remain valid.\par
\noindent\textbf{Later retrieval:} Explain what transferred, what required help and which prerequisite needs another small experiment.\par
\noindent\textbf{Alignment references:} portswigger-mystery, kettle-research, owasp-asvs\par
\subsubsection{ai-application-security: foundations}
\noindent\textbf{Milestone:} foundations-milestone\par
\noindent\textbf{Skill:} Prove the entry skills in a small lab\par
\noindent\textbf{Experiment:} Demonstrate each entry skill with a small owned fixture or permitted training environment. Record the actual setup, one observation per skill and any prerequisite you need to repair.\par
\noindent\textbf{Transfer:} Change a relevant identity, input or environment setting and predict the effect before repeating the demonstration.\par
\noindent\textbf{Control:} Include a deliberately unmet prerequisite and distinguish setup failure from a result about the subject.\par
\noindent\textbf{Later retrieval:} Reconstruct the environment and entry assumptions without following the setup notes.\par
\noindent\textbf{Alignment references:} owasp-llm, owasp-asvs\par
\subsubsection{ai-application-security: boundaries}
\noindent\textbf{Milestone:} boundaries-milestone\par
\noindent\textbf{Skill:} Treat generated text as untrusted input\par
\noindent\textbf{Experiment:} Build a threat model that identifies which component may authorize each action and where untrusted text enters.\par
\noindent\textbf{Transfer:} Add an external document source and revise the trust assumptions.\par
\noindent\textbf{Control:} Reject the idea that a model's assertion of permission is an authorization decision.\par
\noindent\textbf{Later retrieval:} Reconstruct the tool-action boundary without referring to the system prompt.\par
\noindent\textbf{Alignment references:} owasp-llm, owasp-asvs\par
\subsubsection{ai-application-security: retrieval}
\noindent\textbf{Milestone:} retrieval-milestone\par
\noindent\textbf{Skill:} Keep a document from becoming an instruction\par
\noindent\textbf{Experiment:} In a deterministic toy pipeline, test adversarial documents and enforce authorization outside the mock model before any data access or action.\par
\noindent\textbf{Transfer:} Change the document wording or source label and verify the same application boundary still holds.\par
\noindent\textbf{Control:} Keep a legitimate document answer working and deny a requested unrelated data access.\par
\noindent\textbf{Later retrieval:} Explain why prompt wording alone cannot enforce a data-access policy.\par
\noindent\textbf{Alignment references:} owasp-llm, owasp-asvs\par
\subsubsection{ai-application-security: output}
\noindent\textbf{Milestone:} output-milestone\par
\noindent\textbf{Skill:} Validate what the application will execute\par
\noindent\textbf{Experiment:} Add strict schemas and safe handling to a toy assistant's output pipeline and test malformed or unauthorized synthetic outputs.\par
\noindent\textbf{Transfer:} Nest a field or change output encoding and predict whether validation rejects it.\par
\noindent\textbf{Control:} Preserve a valid allowed output while rejecting an extra privileged field.\par
\noindent\textbf{Later retrieval:} Explain schema validity versus authorization of a valid action.\par
\noindent\textbf{Alignment references:} owasp-llm, owasp-asvs\par
\subsubsection{ai-application-security: tools}
\noindent\textbf{Milestone:} tools-milestone\par
\noindent\textbf{Skill:} Give an agent a bounded capability\par
\noindent\textbf{Experiment:} Implement or test an allowlisted dispatcher with synthetic files and a mock model; verify that model output cannot expand the allowed resource set.\par
\noindent\textbf{Transfer:} Change a resource identifier or attempt a second action outside the allowlist.\par
\noindent\textbf{Control:} Permit the intended action and reject traversal-like or unrelated resource references.\par
\noindent\textbf{Later retrieval:} Reconstruct where least privilege is enforced independently of the model.\par
\noindent\textbf{Alignment references:} owasp-llm, owasp-asvs\par
\subsubsection{ai-application-security: data}
\noindent\textbf{Milestone:} data-milestone\par
\noindent\textbf{Skill:} Test retrieval separation between users\par
\noindent\textbf{Experiment:} Test retrieval filters and result handling in an owned fixture so one user's documents cannot be returned to another user.\par
\noindent\textbf{Transfer:} Change a document ID, cache key or role and independently retest isolation.\par
\noindent\textbf{Control:} Include a shared public document that remains legitimately accessible to both users.\par
\noindent\textbf{Later retrieval:} Explain why hiding a private citation after retrieval is too late to guarantee isolation.\par
\noindent\textbf{Alignment references:} owasp-llm, owasp-asvs\par
\subsubsection{ai-application-security: evaluation}
\noindent\textbf{Milestone:} evaluation-milestone\par
\noindent\textbf{Skill:} State exactly what the evaluation measured\par
\noindent\textbf{Experiment:} Run the cases against the toy system, record failures and repairs, and distinguish deterministic application tests from any real-model observations.\par
\noindent\textbf{Transfer:} Add an unseen wording or tool-output shape and rerun the relevant boundary checks.\par
\noindent\textbf{Control:} Include benign lookalikes so blocking every input cannot pass the evaluation.\par
\noindent\textbf{Later retrieval:} Explain why a passing finite test set does not prove universal prompt-injection resistance.\par
\noindent\textbf{Alignment references:} owasp-llm, owasp-asvs\par
\subsubsection{ai-application-security: independent-transfer}
\noindent\textbf{Milestone:} independent-transfer-milestone\par
\noindent\textbf{Skill:} Revisit the skill on a fresh problem\par
\noindent\textbf{Experiment:} After a delay, attempt the fresh task and produce a reproducible evidence packet. Ask a peer to select or inspect the variation when possible; record assistance, failed hypotheses and remaining limits.\par
\noindent\textbf{Transfer:} Change a mechanism-relevant condition, rather than only resetting the old environment, and independently predict and test the resulting behavior.\par
\noindent\textbf{Control:} Include a contrasting case where the mechanism should not apply and a legitimate operation that should remain valid.\par
\noindent\textbf{Later retrieval:} Explain what transferred, what required help and which prerequisite needs another small experiment.\par
\noindent\textbf{Alignment references:} owasp-llm, owasp-asvs\par

\subsection{Research evidence register}
Reviewed 2026-11-01. Titles identify sources; findings below are original paraphrases. Public pages and abstracts are distinguished from full-course or experimental validation.
\subsubsection{retrieval-2006}
\noindent\textbf{Title:} Test-enhanced learning: taking memory tests improves long-term retention - PubMed\par
\noindent\textbf{Evidence:} experimental research; abstract reviewed\par
\noindent\textbf{Supported inference:} Retrieval outperformed restudy on delayed prose-retention tests in these experiments.\par
\noindent\textbf{Limits:} Not an experiment on cybersecurity training or on Ascend.\par
\noindent Access: 2026-11-01, HTTP 200. \url{https://pubmed.ncbi.nlm.nih.gov/16507066/}
\subsubsection{retrieval-2011}
\noindent\textbf{Title:} Retrieval practice produces more learning than elaborative studying with concept mapping - PubMed\par
\noindent\textbf{Evidence:} experimental research; abstract reviewed\par
\noindent\textbf{Supported inference:} Retrieval practice improved comprehension/inference on the science-text assessments studied.\par
\noindent\textbf{Limits:} Does not establish transfer to operational security work or prove a universal superiority for every task.\par
\noindent Access: 2026-11-01, HTTP 200. \url{https://pubmed.ncbi.nlm.nih.gov/21252317/}
\subsubsection{dunlosky-2013}
\noindent\textbf{Title:} Improving Students’ Learning With Effective Learning Techniques: Promising Directions From Cognitive and Educational Psychology – Association for Psychological Science – APS\par
\noindent\textbf{Evidence:} research review; publisher summary reviewed\par
\noindent\textbf{Supported inference:} Practice testing and distributed practice receive stronger support than passive rereading in the review.\par
\noindent\textbf{Limits:} The full underlying set of studies was not independently replicated or audited here.\par
\noindent Access: 2026-11-01, HTTP 200. \url{https://www.psychologicalscience.org/publications/journals/pspi/learning-techniques.html}
\subsubsection{ies-practice}
\noindent\textbf{Title:} WWC | Organizing Instruction and Study to Improve Student Learning\par
\noindent\textbf{Evidence:} evidence-based teaching practice guide\par
\noindent\textbf{Supported inference:} Space study; alternate worked examples and problem solving; connect concrete and abstract representations; ask explanatory questions.\par
\noindent\textbf{Limits:} Recommendations span educational content learning. Ascend’s chosen 2–7 day interval is a pragmatic suggestion, not a measured optimum.\par
\noindent Access: 2026-11-01, HTTP 200. \url{https://ies.ed.gov/ncee/wwc/PracticeGuide/1}
\subsubsection{htb-foundations}
\noindent\textbf{Title:} Information Security Foundations Skill Paths | HTB Academy\par
\noindent\textbf{Evidence:} published curriculum outline\par
\noindent\textbf{Supported inference:} Linux, Windows, scripting, networking and web fundamentals precede specialized security work.\par
\noindent\textbf{Limits:} Public outline reviewed; paid modules and learner outcomes were not audited.\par
\noindent Access: 2026-11-01, HTTP 200. \url{https://academy.hackthebox.com/path/preview/information-security-foundations}
\subsubsection{htb-pentest}
\noindent\textbf{Title:} Penetration Tester Job Role Path | HTB Academy\par
\noindent\textbf{Evidence:} published curriculum outline\par
\noindent\textbf{Supported inference:} Combines foundations, process, enumeration, services, platform-specific escalation, and reporting.\par
\noindent\textbf{Limits:} Ascend covers a subset and is not equivalent to the full job-role path or a professional certification.\par
\noindent Access: 2026-11-01, HTTP 200. \url{https://academy.hackthebox.com/path/preview/penetration-tester}
\subsubsection{htb-soc}
\noindent\textbf{Title:} SOC Analyst Job Role Path | HTB Academy\par
\noindent\textbf{Evidence:} published curriculum outline\par
\noindent\textbf{Supported inference:} Includes incident handling, telemetry, SIEM query work, investigation, and reporting.\par
\noindent\textbf{Limits:} Ascend is narrower, especially in Windows/enterprise DFIR and threat intelligence.\par
\noindent Access: 2026-11-01, HTTP 200. \url{https://academy.hackthebox.com/path/preview/soc-analyst}
\subsubsection{portswigger-paths}
\noindent\textbf{Title:} Learning paths | Web Security Academy - PortSwigger\par
\noindent\textbf{Evidence:} published curriculum landing page\par
\noindent\textbf{Supported inference:} Provides curated web security pathways and links to topic labs.\par
\noindent\textbf{Limits:} Public page reviewed; no authenticated progress or successful lab execution is claimed.\par
\noindent Access: 2026-11-01, HTTP 200. \url{https://portswigger.net/web-security/learning-paths}
\subsubsection{portswigger-mystery}
\noindent\textbf{Title:} Mystery lab challenge | Web Security Academy - PortSwigger\par
\noindent\textbf{Evidence:} published assessment mechanism\par
\noindent\textbf{Supported inference:} Hiding the lab title/category requires discovery and diagnosis instead of matching a supplied vulnerability label.\par
\noindent\textbf{Limits:} Used as an assessment-design reference, not as a guarantee that a beginner can solve a random advanced lab.\par
\noindent Access: 2026-11-01, HTTP 200. \url{https://portswigger.net/web-security/mystery-lab-challenge}
\subsubsection{pwn-college}
\noindent\textbf{Title:} pwn.college\par
\noindent\textbf{Evidence:} university-associated practical curriculum outline\par
\noindent\textbf{Supported inference:} Public modules include web, communication, cryptography, access control, reverse engineering, binary exploitation, and integrated security.\par
\noindent\textbf{Limits:} Public challenges/outline inspected without logging in or solving them. No endorsement by ASU is implied.\par
\noindent Access: 2026-11-01, HTTP 200. \url{https://pwn.college/intro-to-cybersecurity/}
\subsubsection{rop-guide}
\noindent\textbf{Title:} Beginners' guide\par
\noindent\textbf{Evidence:} official challenge pedagogy and technical guide\par
\noindent\textbf{Supported inference:} Introduces isolated ROP concepts with increasing complexity and explicitly distinguishes architectures and calling conventions.\par
\noindent\textbf{Limits:} ROP-only challenge performance is not evidence of broad bug-discovery skill.\par
\noindent Access: 2026-11-01, HTTP 200. \url{https://ropemporium.com/guide.html}
\subsubsection{phoenix-stack0}
\noindent\textbf{Title:} Stack Zero :: Andrew Griffiths' Exploit Education\par
\noindent\textbf{Evidence:} official exercise specification\par
\noindent\textbf{Supported inference:} Stack Zero changes adjacent data; it is not by itself a saved-return-address or code-execution exercise.\par
\noindent\textbf{Limits:} Specification inspected; the challenge was not executed in this audit.\par
\noindent Access: 2026-11-01, HTTP 200. \url{https://exploit.education/phoenix/stack-zero/}
\subsubsection{security-onion-hardware}
\noindent\textbf{Title:} Hardware Requirements — Security Onion Documentation 2.4 documentation\par
\noindent\textbf{Evidence:} official operational requirements\par
\noindent\textbf{Supported inference:} Security Onion supports x86-64 only; Import, Eval and Standalone have different resource requirements.\par
\noindent\textbf{Limits:} The Import table says 50 GB while the first-time guide says 200 GB. Plan conservatively for 200 GB and check the chosen release.\par
\noindent Access: 2026-11-01, HTTP 200. \url{https://docs.securityonion.net/en/2.4/hardware.html}
\subsubsection{security-onion-import}
\noindent\textbf{Title:} First Time Users — Security Onion Documentation 2.4 documentation\par
\noindent\textbf{Evidence:} official workflow guide\par
\noindent\textbf{Supported inference:} Import accepts PCAP/EVTX datasets for investigation; it is not a full continuous-monitoring deployment.\par
\noindent\textbf{Limits:} Minimum storage wording differs from the hardware table; no actual installation was performed.\par
\noindent Access: 2026-11-01, HTTP 200. \url{https://docs.securityonion.net/en/2.4/first-time-users.html}
\subsubsection{evans-learning}
\noindent\textbf{Title:} How to teach yourself hard things\par
\noindent\textbf{Evidence:} practitioner first-person learning account\par
\noindent\textbf{Supported inference:} Julia Evans describes turning vague confusion into specific questions, investigating them, and calibrating understanding.\par
\noindent\textbf{Limits:} A reflective personal account, not controlled causal evidence or a promise of the same career outcome.\par
\noindent Access: 2026-11-01, HTTP 200. \url{https://jvns.ca/blog/2018/09/01/learning-skills-you-can-practice/}
\subsubsection{liveoverflow-learning}
\noindent\textbf{Title:} How Do I Get Started With Hacking?\par
\noindent\textbf{Evidence:} practitioner first-person advice\par
\noindent\textbf{Supported inference:} Programming, CTF/wargame attempts, and studying missed solutions are recurring practices in the author’s account.\par
\noindent\textbf{Limits:} Do not adopt the page’s 10,000-hour slogan as a universal law or copy its dismissal of all paid courses.\par
\noindent Access: 2026-11-01, HTTP 200. \url{https://liveoverflow.com/faq/get-started-with-hacking/}
\subsubsection{vavra-interview}
\noindent\textbf{Title:} Interview with Kamil Vavra - Hall of Fame high flyer | Web Security Academy - PortSwigger\par
\noindent\textbf{Evidence:} vendor-hosted practitioner interview\par
\noindent\textbf{Supported inference:} Kamil Vavra reports translating lab practice and new vulnerability concepts into workplace findings, and getting help when stuck.\par
\noindent\textbf{Limits:} Vendor selection and survivorship bias; self-reported success is corroborating context, not proof of typical outcomes. All-night or compulsive behavior is not recommended.\par
\noindent Access: 2026-11-01, HTTP 200. \url{https://portswigger.net/web-security/getting-started/kamil-vavra/index.html}
\subsubsection{villarreal-interview}
\noindent\textbf{Title:} Interview with Johnny Villarreal - Hall of Fame high flyer | Web Security Academy - PortSwigger\par
\noindent\textbf{Evidence:} vendor-hosted practitioner interview\par
\noindent\textbf{Supported inference:} Johnny Villarreal reports confidence gains from practical labs and argues for beginner scaffolding around difficult logic puzzles.\par
\noindent\textbf{Limits:} Self-report from a selected successful learner; does not isolate the effect of the platform.\par
\noindent Access: 2026-11-01, HTTP 200. \url{https://portswigger.net/web-security/getting-started/johnny-villarreal/index.html}
\subsubsection{kettle-research}
\noindent\textbf{Title:} Server-Side Template Injection | PortSwigger Research\par
\noindent\textbf{Evidence:} practitioner technical research account\par
\noindent\textbf{Supported inference:} James Kettle’s research demonstrates mechanism-based diagnosis and comparison across distinct implementations.\par
\noindent\textbf{Limits:} Technical research methodology is informative but not evidence that this curriculum produces a comparable researcher.\par
\noindent Access: 2026-11-01, HTTP 200. \url{https://portswigger.net/research/server-side-template-injection}
\subsubsection{pwn-binary}
\noindent\textbf{Title:} pwn.college: Binary Exploitation\par
\noindent\textbf{Evidence:} official public resource; page inspected\par
\noindent\textbf{Supported inference:} The public module separates adjacent data corruption, control hijack, canaries/ASLR, shellcode and integrated exercises.\par
\noindent\textbf{Limits:} Public module text was inspected; challenges were not all executed or assessed.\par
\noindent Access: 2026-11-01, HTTP 200. \url{https://pwn.college/intro-to-cybersecurity/binary-exploitation/}
\subsubsection{phoenix-index}
\noindent\textbf{Title:} Exploit Education: Phoenix index\par
\noindent\textbf{Evidence:} official public resource; page inspected\par
\noindent\textbf{Supported inference:} The official navigation provides separate stack, format-string and heap exercise sequences and setup instructions.\par
\noindent\textbf{Limits:} An index is a route to exercises, not a complete theory text or proof of modern exploit portability.\par
\noindent Access: 2026-11-01, HTTP 200. \url{https://exploit.education/phoenix/}
\subsubsection{rapid7-db}
\noindent\textbf{Title:} Rapid7 Vulnerability and Exploit Database\par
\noindent\textbf{Evidence:} official public resource; page inspected\par
\noindent\textbf{Supported inference:} The public database is an available starting point for matching service findings with documented vulnerability and exploit conditions.\par
\noindent\textbf{Limits:} Individual entries still require version, configuration and source review; a listed exploit is not guaranteed to apply.\par
\noindent Access: 2026-11-01, HTTP 200. \url{https://www.rapid7.com/db/}
\subsubsection{mitre-emulation}
\noindent\textbf{Title:} GitHub - center-for-threat-informed-defense/adversary\_emulation\_library: An open library of adversary emulation plans designed to empower organizations to test their defenses based\par
\noindent\textbf{Evidence:} first-party public curriculum, technical guidance or project documentation; page reviewed\par
\noindent\textbf{Supported inference:} Public behavior-based plans support objective-led emulation and explicit prerequisites.\par
\noindent\textbf{Limits:} Alignment or implementation guidance, not empirical validation of Ascend learning outcomes. Private course contents and every external exercise were not audited.\par
\noindent Access: 2026-11-01, HTTP 200. \url{https://github.com/center-for-threat-informed-defense/adversary_emulation_library}
\subsubsection{offsec-pen}
\noindent\textbf{Title:} Get your OSCP+ certification with PEN-200 | OffSec\par
\noindent\textbf{Evidence:} first-party public curriculum, technical guidance or project documentation; page reviewed\par
\noindent\textbf{Supported inference:} The public outline combines enumeration, web, host and directory testing and evidence gathering.\par
\noindent\textbf{Limits:} Alignment or implementation guidance, not empirical validation of Ascend learning outcomes. Private course contents and every external exercise were not audited.\par
\noindent Access: 2026-11-01, HTTP 200. \url{https://www.offsec.com/courses/pen-200/}
\subsubsection{microsoft-soc}
\noindent\textbf{Title:} Study guide for Exam SC-200: Microsoft Security Operations Analyst | Microsoft Learn\par
\noindent\textbf{Evidence:} first-party public curriculum, technical guidance or project documentation; page reviewed\par
\noindent\textbf{Supported inference:} The role outline separates security operations, incident response and hunting.\par
\noindent\textbf{Limits:} Alignment or implementation guidance, not empirical validation of Ascend learning outcomes. Private course contents and every external exercise were not audited.\par
\noindent Access: 2026-11-01, HTTP 200. \url{https://learn.microsoft.com/en-us/credentials/certifications/resources/study-guides/sc-200}
\subsubsection{microsoft-ad}
\noindent\textbf{Title:} Active Directory Domain Services overview | Microsoft Learn\par
\noindent\textbf{Evidence:} first-party public curriculum, technical guidance or project documentation; page reviewed\par
\noindent\textbf{Supported inference:} Directory objects, discovery, authentication and permissions are distinct AD DS responsibilities.\par
\noindent\textbf{Limits:} Alignment or implementation guidance, not empirical validation of Ascend learning outcomes. Private course contents and every external exercise were not audited.\par
\noindent Access: 2026-11-01, HTTP 200. \url{https://learn.microsoft.com/en-us/windows-server/identity/ad-ds/get-started/virtual-dc/active-directory-domain-services-overview}
\subsubsection{microsoft-adcs}
\noindent\textbf{Title:} What is Active Directory Certificate Services in Windows Server? | Microsoft Learn\par
\noindent\textbf{Evidence:} first-party public curriculum, technical guidance or project documentation; page reviewed\par
\noindent\textbf{Supported inference:} AD CS is a Windows Server PKI role with certificate enrollment and identity-use policies.\par
\noindent\textbf{Limits:} Alignment or implementation guidance, not empirical validation of Ascend learning outcomes. Private course contents and every external exercise were not audited.\par
\noindent Access: 2026-11-01, HTTP 200. \url{https://learn.microsoft.com/en-us/windows-server/identity/ad-cs/active-directory-certificate-services-overview}
\subsubsection{owasp-cloud}
\noindent\textbf{Title:} OWASP Foundation\par
\noindent\textbf{Evidence:} first-party public curriculum, technical guidance or project documentation; page reviewed\par
\noindent\textbf{Supported inference:} The interim risk list covers configuration, identity, supply chain, secrets and network boundaries.\par
\noindent\textbf{Limits:} Alignment or implementation guidance, not empirical validation of Ascend learning outcomes. Private course contents and every external exercise were not audited.\par
\noindent Access: 2026-11-01, HTTP 200. \url{https://owasp.org/www-project-cloud-native-application-security-top-10/}
\subsubsection{kubernetes-checklist}
\noindent\textbf{Title:} Application Security Checklist | Kubernetes\par
\noindent\textbf{Evidence:} first-party public curriculum, technical guidance or project documentation; page reviewed\par
\noindent\textbf{Supported inference:} The official checklist covers service accounts, security contexts, RBAC, images and network policies.\par
\noindent\textbf{Limits:} Alignment or implementation guidance, not empirical validation of Ascend learning outcomes. Private course contents and every external exercise were not audited.\par
\noindent Access: 2026-11-01, HTTP 200. \url{https://kubernetes.io/docs/concepts/security/application-security-checklist/}
\subsubsection{owasp-asvs}
\noindent\textbf{Title:} Overview | OWASP Application Security Verification Standard\par
\noindent\textbf{Evidence:} first-party public curriculum, technical guidance or project documentation; page reviewed\par
\noindent\textbf{Supported inference:} ASVS supplies testable application security requirements, not a certification of an untested implementation.\par
\noindent\textbf{Limits:} Alignment or implementation guidance, not empirical validation of Ascend learning outcomes. Private course contents and every external exercise were not audited.\par
\noindent Access: 2026-11-01, HTTP 200. \url{https://owasp.org/www-project-application-security-verification-standard/}
\subsubsection{digitalcorpora}
\noindent\textbf{Title:} Scenarios – Digital Corpora\par
\noindent\textbf{Evidence:} first-party public curriculum, technical guidance or project documentation; page reviewed\par
\noindent\textbf{Supported inference:} The public scenario index provides training evidence; dataset terms and storage requirements still need checking.\par
\noindent\textbf{Limits:} Alignment or implementation guidance, not empirical validation of Ascend learning outcomes. Private course contents and every external exercise were not audited.\par
\noindent Access: 2026-11-01, HTTP 200. \url{https://digitalcorpora.org/corpora/scenarios/}
\subsubsection{stix}
\noindent\textbf{Title:} Introduction to STIX\par
\noindent\textbf{Evidence:} first-party public curriculum, technical guidance or project documentation; page reviewed\par
\noindent\textbf{Supported inference:} Typed objects and relationships convey threat information; syntax validation cannot establish the underlying claim.\par
\noindent\textbf{Limits:} Alignment or implementation guidance, not empirical validation of Ascend learning outcomes. Private course contents and every external exercise were not audited.\par
\noindent Access: 2026-11-01, HTTP 200. \url{https://oasis-open.github.io/cti-documentation/stix/intro.html}
\subsubsection{flare-vm}
\noindent\textbf{Title:} GitHub - mandiant/flare-vm: A collection of software installations scripts for Windows systems that allows you to easily setup and maintain a reverse engineering environment on a V\par
\noindent\textbf{Evidence:} first-party public curriculum, technical guidance or project documentation; page reviewed\par
\noindent\textbf{Supported inference:} The analysis toolkit is designed for an appropriate Windows virtual environment.\par
\noindent\textbf{Limits:} Alignment or implementation guidance, not empirical validation of Ascend learning outcomes. Private course contents and every external exercise were not audited.\par
\noindent Access: 2026-11-01, HTTP 200. \url{https://github.com/mandiant/flare-vm}
\subsubsection{remnux}
\noindent\textbf{Title:} Get the Virtual Appliance | REMnux Documentation\par
\noindent\textbf{Evidence:} first-party public curriculum, technical guidance or project documentation; page reviewed\par
\noindent\textbf{Supported inference:} The current guide offers both AMD64 and ARM64 appliances with hypervisor-specific instructions.\par
\noindent\textbf{Limits:} Alignment or implementation guidance, not empirical validation of Ascend learning outcomes. Private course contents and every external exercise were not audited.\par
\noindent Access: 2026-11-01, HTTP 200. \url{https://docs.remnux.org/install-distro/get-virtual-appliance}
\subsubsection{mastg}
\noindent\textbf{Title:} OWASP MASTG - OWASP Mobile Application Security\par
\noindent\textbf{Evidence:} first-party public curriculum, technical guidance or project documentation; page reviewed\par
\noindent\textbf{Supported inference:} Mobile tests cover platform behavior, storage, transport, components and application code.\par
\noindent\textbf{Limits:} Alignment or implementation guidance, not empirical validation of Ascend learning outcomes. Private course contents and every external exercise were not audited.\par
\noindent Access: 2026-11-01, HTTP 200. \url{https://mas.owasp.org/MASTG/}
\subsubsection{owasp-api}
\noindent\textbf{Title:} OWASP Top 10 API Security Risks – 2023 - OWASP API Security Top 10\par
\noindent\textbf{Evidence:} first-party public curriculum, technical guidance or project documentation; page reviewed\par
\noindent\textbf{Supported inference:} The 2023 risks distinguish object/property/function authorization, resource consumption and workflow risks.\par
\noindent\textbf{Limits:} Alignment or implementation guidance, not empirical validation of Ascend learning outcomes. Private course contents and every external exercise were not audited.\par
\noindent Access: 2026-11-01, HTTP 200. \url{https://owasp.org/API-Security/editions/2023/en/0x11-t10/}
\subsubsection{cryptohack}
\noindent\textbf{Title:} CryptoHack – Courses\par
\noindent\textbf{Evidence:} first-party public curriculum, technical guidance or project documentation; page reviewed\par
\noindent\textbf{Supported inference:} Public courses progress from representations and arithmetic to symmetric and public-key mechanisms.\par
\noindent\textbf{Limits:} Alignment or implementation guidance, not empirical validation of Ascend learning outcomes. Private course contents and every external exercise were not audited.\par
\noindent Access: 2026-11-01, HTTP 200. \url{https://cryptohack.org/courses/}
\subsubsection{owasp-iot}
\noindent\textbf{Title:} OWASP Internet of Things | OWASP Foundation\par
\noindent\textbf{Evidence:} first-party public curriculum, technical guidance or project documentation; page reviewed\par
\noindent\textbf{Supported inference:} The project covers device, interface, data, firmware-update and ecosystem security concerns.\par
\noindent\textbf{Limits:} Alignment or implementation guidance, not empirical validation of Ascend learning outcomes. Private course contents and every external exercise were not audited.\par
\noindent Access: 2026-11-01, HTTP 200. \url{https://owasp.org/www-project-internet-of-things/}
\subsubsection{mosquitto}
\noindent\textbf{Title:} Documentation | Eclipse Mosquitto\par
\noindent\textbf{Evidence:} first-party public curriculum, technical guidance or project documentation; page reviewed\par
\noindent\textbf{Supported inference:} Official references describe broker/client behavior, configuration, authentication and topic controls.\par
\noindent\textbf{Limits:} Alignment or implementation guidance, not empirical validation of Ascend learning outcomes. Private course contents and every external exercise were not audited.\par
\noindent Access: 2026-11-01, HTTP 200. \url{https://mosquitto.org/documentation/}
\subsubsection{cisa-ics}
\noindent\textbf{Title:} ICS Recommended Practices | CISA\par
\noindent\textbf{Evidence:} first-party public curriculum, technical guidance or project documentation; page reviewed\par
\noindent\textbf{Supported inference:} Industrial guidance addresses operational dependencies and safety as well as network protection.\par
\noindent\textbf{Limits:} Alignment or implementation guidance, not empirical validation of Ascend learning outcomes. Private course contents and every external exercise were not audited.\par
\noindent Access: 2026-11-01, HTTP 200. \url{https://www.cisa.gov/resources-tools/resources/ics-recommended-practices}
\subsubsection{pymodbus}
\noindent\textbf{Title:} Welcome to PyModbus’s documentation! — PyModbus 4.0.0dev16 documentation\par
\noindent\textbf{Evidence:} first-party public curriculum, technical guidance or project documentation; page reviewed\par
\noindent\textbf{Supported inference:} The project provides clients, servers and simulators with version-specific APIs.\par
\noindent\textbf{Limits:} Alignment or implementation guidance, not empirical validation of Ascend learning outcomes. Private course contents and every external exercise were not audited.\par
\noindent Access: 2026-11-01, HTTP 200. \url{https://pymodbus.readthedocs.io/en/latest/}
\subsubsection{nist-csf}
\noindent\textbf{Title:} Cybersecurity Framework | NIST\par
\noindent\textbf{Evidence:} first-party public curriculum, technical guidance or project documentation; page reviewed\par
\noindent\textbf{Supported inference:} The framework supplies outcomes for governance and cybersecurity risk management.\par
\noindent\textbf{Limits:} Alignment or implementation guidance, not empirical validation of Ascend learning outcomes. Private course contents and every external exercise were not audited.\par
\noindent Access: 2026-11-01, HTTP 200. \url{https://www.nist.gov/cyberframework}
\subsubsection{cisa-kev}
\noindent\textbf{Title:} Known Exploited Vulnerabilities Catalog | CISA\par
\noindent\textbf{Evidence:} first-party public curriculum, technical guidance or project documentation; page reviewed\par
\noindent\textbf{Supported inference:} Known exploitation informs prioritization; catalog membership does not prove compromise of a particular system.\par
\noindent\textbf{Limits:} Alignment or implementation guidance, not empirical validation of Ascend learning outcomes. Private course contents and every external exercise were not audited.\par
\noindent Access: 2026-11-01, HTTP 200. \url{https://www.cisa.gov/known-exploited-vulnerabilities-catalog}
\subsubsection{epss}
\noindent\textbf{Title:} Exploit Prediction Scoring System (EPSS)\par
\noindent\textbf{Evidence:} first-party public curriculum, technical guidance or project documentation; page reviewed\par
\noindent\textbf{Supported inference:} EPSS models exploitation probability, distinct from severity and an organization's specific exposure.\par
\noindent\textbf{Limits:} Alignment or implementation guidance, not empirical validation of Ascend learning outcomes. Private course contents and every external exercise were not audited.\par
\noindent Access: 2026-11-01, HTTP 200. \url{https://www.first.org/epss/}
\subsubsection{cisa-ssvc}
\noindent\textbf{Title:} Stakeholder-Specific Vulnerability Categorization (SSVC) | CISA\par
\noindent\textbf{Evidence:} first-party public curriculum, technical guidance or project documentation; page reviewed\par
\noindent\textbf{Supported inference:} SSVC uses stakeholder context and decision paths to prioritize vulnerability response.\par
\noindent\textbf{Limits:} Alignment or implementation guidance, not empirical validation of Ascend learning outcomes. Private course contents and every external exercise were not audited.\par
\noindent Access: 2026-11-01, HTTP 200. \url{https://www.cisa.gov/stakeholder-specific-vulnerability-categorization-ssvc}
\subsubsection{owasp-llm}
\noindent\textbf{Title:} LLMRisks Archive - OWASP Gen AI Security Project\par
\noindent\textbf{Evidence:} first-party public curriculum, technical guidance or project documentation; page reviewed\par
\noindent\textbf{Supported inference:} The 2025 risks include prompt injection, data boundaries, output handling, excessive agency and resource use.\par
\noindent\textbf{Limits:} Alignment or implementation guidance, not empirical validation of Ascend learning outcomes. Private course contents and every external exercise were not audited.\par
\noindent Access: 2026-11-01, HTTP 200. \url{https://genai.owasp.org/llm-top-10/}
\subsubsection{juice-docker}
\noindent\textbf{Title:} Running OWASP Juice Shop :: Pwning OWASP Juice Shop\par
\noindent\textbf{Evidence:} first-party public curriculum, technical guidance or project documentation; page reviewed\par
\noindent\textbf{Supported inference:} The guide documents Docker and ARM variants; actual current image manifests were checked separately.\par
\noindent\textbf{Limits:} Alignment or implementation guidance, not empirical validation of Ascend learning outcomes. Private course contents and every external exercise were not audited.\par
\noindent Access: 2026-11-01, HTTP 200. \url{https://pwning.owasp-juice.shop/companion-guide/latest/part1/running.html}
\subsubsection{juice-repo}
\noindent\textbf{Title:} GitHub - juice-shop/juice-shop: OWASP Juice Shop: Probably the most modern and sophisticated insecure web application · GitHub\par
\noindent\textbf{Evidence:} first-party public curriculum, technical guidance or project documentation; page reviewed\par
\noindent\textbf{Supported inference:} The intentionally vulnerable application records challenge completion for its running instance.\par
\noindent\textbf{Limits:} Alignment or implementation guidance, not empirical validation of Ascend learning outcomes. Private course contents and every external exercise were not audited.\par
\noindent Access: 2026-11-01, HTTP 200. \url{https://github.com/juice-shop/juice-shop}
\subsubsection{webgoat}
\noindent\textbf{Title:} GitHub - WebGoat/WebGoat: WebGoat is a deliberately insecure application · GitHub\par
\noindent\textbf{Evidence:} first-party public curriculum, technical guidance or project documentation; page reviewed\par
\noindent\textbf{Supported inference:} Interactive lessons are designed for controlled training environments with limited network exposure.\par
\noindent\textbf{Limits:} Alignment or implementation guidance, not empirical validation of Ascend learning outcomes. Private course contents and every external exercise were not audited.\par
\noindent Access: 2026-11-01, HTTP 200. \url{https://github.com/WebGoat/WebGoat}
\subsubsection{dvwa}
\noindent\textbf{Title:} GitHub - digininja/DVWA: Damn Vulnerable Web Application (DVWA) · GitHub\par
\noindent\textbf{Evidence:} first-party public curriculum, technical guidance or project documentation; page reviewed\par
\noindent\textbf{Supported inference:} The repository provides vulnerable web exercises and deployment instructions under an open license.\par
\noindent\textbf{Limits:} Alignment or implementation guidance, not empirical validation of Ascend learning outcomes. Private course contents and every external exercise were not audited.\par
\noindent Access: 2026-11-01, HTTP 200. \url{https://github.com/digininja/DVWA}
\subsubsection{crapi}
\noindent\textbf{Title:} GitHub - OWASP/crAPI: completely ridiculous API (crAPI) · GitHub\par
\noindent\textbf{Evidence:} first-party public curriculum, technical guidance or project documentation; page reviewed\par
\noindent\textbf{Supported inference:} The deliberately vulnerable API application has a multi-service deployment with separate runtime requirements.\par
\noindent\textbf{Limits:} Alignment or implementation guidance, not empirical validation of Ascend learning outcomes. Private course contents and every external exercise were not audited.\par
\noindent Access: 2026-11-01, HTTP 200. \url{https://github.com/OWASP/crAPI}
\subsubsection{goad}
\noindent\textbf{Title:} GitHub - Orange-Cyberdefense/GOAD: game of active directory · GitHub\par
\noindent\textbf{Evidence:} first-party public curriculum, technical guidance or project documentation; page reviewed\par
\noindent\textbf{Supported inference:} GOAD provisions Windows VMs; its Docker launcher does not turn the domain controllers into Linux containers.\par
\noindent\textbf{Limits:} Alignment or implementation guidance, not empirical validation of Ascend learning outcomes. Private course contents and every external exercise were not audited.\par
\noindent Access: 2026-11-01, HTTP 200. \url{https://github.com/Orange-Cyberdefense/GOAD}
\subsubsection{go-labs}
\noindent\textbf{Title:} GitHub - dolevf/Damn-Vulnerable-GraphQL-Application: Damn Vulnerable GraphQL Application is an intentionally vulnerable GraphQL service implementation designed for learning about a\par
\noindent\textbf{Evidence:} first-party public curriculum, technical guidance or project documentation; page reviewed\par
\noindent\textbf{Supported inference:} The repository offers a vulnerable GraphQL training application; current architecture support needs separate verification.\par
\noindent\textbf{Limits:} Alignment or implementation guidance, not empirical validation of Ascend learning outcomes. Private course contents and every external exercise were not audited.\par
\noindent Access: 2026-11-01, HTTP 200. \url{https://github.com/dolevf/Damn-Vulnerable-GraphQL-Application}
\subsubsection{atomic-red}
\noindent\textbf{Title:} GitHub - redcanaryco/atomic-red-team: Small and highly portable detection tests based on MITRE's ATT\&CK. · GitHub\par
\noindent\textbf{Evidence:} first-party public curriculum, technical guidance or project documentation; page reviewed\par
\noindent\textbf{Supported inference:} Individual emulation tests require review of commands, prerequisites and cleanup before execution.\par
\noindent\textbf{Limits:} Alignment or implementation guidance, not empirical validation of Ascend learning outcomes. Private course contents and every external exercise were not audited.\par
\noindent Access: 2026-11-01, HTTP 200. \url{https://github.com/redcanaryco/atomic-red-team}
\subsubsection{cloudgoat}
\noindent\textbf{Title:} GitHub - RhinoSecurityLabs/cloudgoat: CloudGoat is Rhino Security Labs' "Vulnerable by Design" AWS deployment tool · GitHub\par
\noindent\textbf{Evidence:} first-party public curriculum, technical guidance or project documentation; page reviewed\par
\noindent\textbf{Supported inference:} The tool provisions real AWS training resources and requires a scoped account, budget and teardown.\par
\noindent\textbf{Limits:} Alignment or implementation guidance, not empirical validation of Ascend learning outcomes. Private course contents and every external exercise were not audited.\par
\noindent Access: 2026-11-01, HTTP 200. \url{https://github.com/RhinoSecurityLabs/cloudgoat}
\subsubsection{python-http-lab}
\noindent\textbf{Title:} Python HTTP server command-line interface\par
\noindent\textbf{Evidence:} technical documentation for local training recipes\par
\noindent\textbf{Supported inference:} The guided lab uses an explicit loopback bind and a disposable document directory.\par
\noindent\textbf{Limits:} Used for command/format semantics; local fixture tests do not establish operational proficiency or measured learning outcomes.\par
\noindent Access: 2026-11-02, HTTP 200. \url{https://docs.python.org/3/library/http.server.html#command-line-interface}
\subsubsection{python-datetime-lab}
\noindent\textbf{Title:} Python datetime reference\par
\noindent\textbf{Evidence:} technical documentation for local training recipes\par
\noindent\textbf{Supported inference:} Timezone-aware ISO timestamps can be normalized to UTC before comparing a log window.\par
\noindent\textbf{Limits:} Used for command/format semantics; local fixture tests do not establish operational proficiency or measured learning outcomes.\par
\noindent Access: 2026-11-02, HTTP 200. \url{https://docs.python.org/3/library/datetime.html}
\subsubsection{wireshark-fixtures}
\noindent\textbf{Title:} Wireshark User Guide\par
\noindent\textbf{Evidence:} technical documentation for local training recipes\par
\noindent\textbf{Supported inference:} Packet inspection distinguishes Ethernet HTTP exchanges from a raw IEEE 802.11 beacon.\par
\noindent\textbf{Limits:} Used for command/format semantics; local fixture tests do not establish operational proficiency or measured learning outcomes.\par
\noindent Access: 2026-11-02, HTTP 200. \url{https://www.wireshark.org/docs/wsug_html_chunked/}
\subsubsection{Discovery limits}
200; title visible, syllabus requires client rendering; not used for detailed curriculum claims. \url{https://www.netacad.com/courses/networking-basics?courseLang=en-US}\par
404; not used as evidence. \url{https://pwn.college/teaching/}\par
Expansion audit returned HTTP 404; no syllabus claim is based on this response. \url{https://academy.hackthebox.com/path/preview/red-teamer}\par
Expansion audit returned HTTP 403; no syllabus claim is based on this response. \url{https://training.zeropointsecurity.co.uk/courses/red-team-ops}\par

